% Generated by build_notes.py from modules/*.md. Edit those sources.
\documentclass[11pt,oneside]{report}
\usepackage[letterpaper,margin=1in,headheight=15pt]{geometry}
\usepackage{fontspec}
\setmainfont{texgyrepagella-regular.otf}[BoldFont=texgyrepagella-bold.otf,ItalicFont=texgyrepagella-italic.otf,BoldItalicFont=texgyrepagella-bolditalic.otf]
\setsansfont{texgyreheros-regular.otf}[BoldFont=texgyreheros-bold.otf,ItalicFont=texgyreheros-italic.otf]
\setmonofont{lmmono10-regular.otf}[Scale=0.85]
\usepackage{amsmath,amssymb,amsthm,booktabs,longtable,array,multirow,calc}
\usepackage{graphicx,xcolor,tikz,pgfplots}
\pgfplotsset{compat=1.18}
\usetikzlibrary{arrows.meta,positioning,calc,fit,shapes.geometric,patterns,decorations.pathreplacing}
% UIUC official brand colors, verified 2026-10-23.
% https://brand.illinois.edu/visual-identity/color
\definecolor{IlliniOrange}{HTML}{FF5F05}
\definecolor{IlliniBlue}{HTML}{13294B}
\definecolor{Storm}{HTML}{707372}
\definecolor{LightStorm}{HTML}{C6C7C6}
\definecolor{Industrial}{HTML}{1D58A7}
\definecolor{Patina}{HTML}{007E8E}
\tikzset{>=Stealth, every node/.style={font=\small,text=IlliniBlue}, box/.style={draw=IlliniBlue,rounded corners=2pt,align=center,inner sep=6pt,execute at begin node=\hyphenpenalty10000\relax}, orangebox/.style={box,draw=IlliniOrange,very thick}, edge/.style={->,thick,IlliniBlue}, accent/.style={->,very thick,IlliniOrange}}
\pgfplotsset{every axis/.append style={width=13.5cm,height=6.5cm,axis line style={IlliniBlue},tick label style={font=\small,text=IlliniBlue},label style={font=\small,text=IlliniBlue},legend style={font=\small,draw=none,fill=white},grid style={LightStorm},major grid style={LightStorm},axis lines=left,clip=false}, every axis plot/.append style={thick}}

\usepackage{titlesec,fancyhdr,enumitem,caption,float,placeins,fvextra}
\usepackage{needspace}
\usepackage{microtype,xurl}
\usepackage{hyperref}
\hypersetup{colorlinks=true,linkcolor=IlliniBlue,urlcolor=Industrial,citecolor=IlliniBlue,pdftitle={BADM 534: Data Analytics for Marketing Applications},pdfauthor={Nathan Yang}}
\usepackage{bookmark}
\titleformat{\chapter}[display]{\normalfont\huge\bfseries\raggedright\hyphenpenalty=10000\color{IlliniBlue}}{\chaptertitlename\ \thechapter}{8pt}{\huge}
\titlespacing*{\chapter}{0pt}{-18pt}{20pt}
\titleformat{\section}{\Large\bfseries\raggedright\hyphenpenalty=10000\color{IlliniBlue}}{\thesection}{0.7em}{}
\titleformat{\subsection}{\large\bfseries\raggedright\hyphenpenalty=10000\color{IlliniBlue}}{\thesubsection}{0.7em}{}
\setlength{\parindent}{1.3em}
\setlength{\parskip}{3pt}
\linespread{1.06}
\setlist{itemsep=3pt,topsep=5pt}
\setcounter{tocdepth}{1}
\setcounter{secnumdepth}{2}
\pagestyle{fancy}
\fancyhf{}
\fancyhead[L]{\small\sffamily BADM 534}
\fancyhead[R]{\small\sffamily Lecture Notes}
\fancyfoot[C]{\thepage}
\renewcommand{\headrulewidth}{0.4pt}
\newcommand{\tightlist}{\setlength{\itemsep}{0pt}\setlength{\parskip}{0pt}}
\DefineVerbatimEnvironment{Highlighting}{Verbatim}{breaklines,fontsize=\small,commandchars=\\\{\}}
\newtheorem{assumption}{Assumption}[chapter]
\captionsetup{font=small,labelfont=bf,labelsep=period}
\setlength{\emergencystretch}{4em}
\providecommand{\pandocbounded}[1]{#1}

\clubpenalty=10000
\widowpenalty=10000

\begin{document}
\hypersetup{pageanchor=false}
\begin{titlepage}
\color{IlliniBlue}
\vspace*{1.2cm}
{\sffamily\Large BADM 534\par}
\vspace{.7cm}
{\Huge\bfseries Data Analytics for\\[.2cm] Marketing Applications\par}
\vspace{.8cm}
{\LARGE Lecture Notes\par}
\vspace{.5cm}
{\large Marketing Decisions and AI-Assisted Coding\par}
\vspace{1.1cm}
\noindent{\color{IlliniOrange}\rule{\textwidth}{2pt}}\par
\vspace{.7cm}
{\large By Nathan Yang\par}
{\large October 23, 2026\par}
\end{titlepage}
\pagenumbering{roman}
\hypersetup{pageanchor=true}
\tableofcontents
\clearpage
\pagenumbering{arabic}
\hypertarget{M01-m1.-marketing-decisions-data-and-ai-assisted-python}{%
\chapter{Marketing Decisions, Data, and AI-Assisted
Python}\label{M01-m1.-marketing-decisions-data-and-ai-assisted-python}}

\noindent A retailer must decide which customers should receive a coupon, how much
inventory to place in each store, and which product descriptions should
appear online. Each decision needs a different analytical object. A
forecast of next week's purchases can inform inventory. A comparison of
purchases with and without a coupon can inform targeting. A summary of
product complaints can identify a service problem. An analysis becomes
useful when its output corresponds to the decision the manager can
actually change.

\hypertarget{M01-from-a-marketing-question-to-an-analytical-object}{%
\section{From a Marketing Question to an Analytical
Object}\label{M01-from-a-marketing-question-to-an-analytical-object}}

\noindent Consider a subscription service deciding whether to offer a renewal
discount. The service observes account tenure, past payments, recent
use, and prior contacts. A descriptive analysis reports renewal rates by
tenure. A predictive analysis estimates the probability that an account
renews under the conditions represented in the training data. A causal
analysis estimates how an offer changes renewal. A decision analysis
combines that change with margins, contact costs, and operational
constraints. These objects can use overlapping data while answering
different questions.

The timing of the decision determines what information is available. If
an account must be selected on Monday morning, a support call made on
Wednesday cannot be a predictor for that selection. A variable can be
strongly associated with an outcome and still be unusable at the
decision time. Write down a forecast origin, an outcome window, and a
feature window before constructing a table. An account-week table and an
account-lifetime table support different forecasts because their rows
and time horizons differ.

For a renewal offer, a useful decision quantity is incremental expected
contribution:

\[
\Delta\pi_i=m_i\{p_i(1)-p_i(0)\}-c_i.
\]

Here \(i\) identifies an account, \(p_i(1)\) and \(p_i(0)\) are its
renewal probabilities with and without the offer, \(m_i\) is
contribution in dollars per retained account over the specified horizon,
and \(c_i\) is the expected incremental cost of making the offer. This
compact expression assumes that the offer's costs can be summarized by
\(c_i\) and that \(m_i\) already reflects relevant variable costs. When
every renewing recipient receives a discount, discount expense depends
on renewal under treatment and should be modeled explicitly.

The difference between the two renewal probabilities matters because
many accounts would renew without assistance. Suppose contribution is
\$100 and contact cost is \$3. Then an account whose renewal probability
rises from 0.90 to 0.91 has incremental expected contribution of \$-2.
An account whose probability rises from 0.40 to 0.50 has contribution of
\$7. Ranking by treated renewal probability chooses the first account;
ranking by incremental contribution chooses the second. Observing only
treated customers would not reveal the untreated probabilities.

\begin{figure}[htbp]
\centering
\begin{tikzpicture}
\begin{axis}[name=left,width=6.3cm,height=6.1cm,xmin=.5,xmax=2.5,ybar,bar width=20pt,ymin=0,ymax=1.05,xtick={1,2},xticklabels={Account A,Account B},ylabel={Renewal probability with offer},ytick={0,0.5,1},xmajorgrids=false,ymajorgrids=true]
\addplot[fill=IlliniBlue,draw=IlliniBlue] coordinates {(1,0.91)(2,0.50)};
\node at (axis cs:1,.98) {0.91};\node at (axis cs:2,.57) {0.50};
\end{axis}
\begin{axis}[at={(left.east)},anchor=west,xshift=2.3cm,width=6.3cm,height=6.1cm,xmin=.5,xmax=2.5,ybar,bar width=20pt,ymin=-4,ymax=9,xtick={1,2},xticklabels={Account A,Account B},ylabel={Incremental contribution (dollars)},ytick={-2,0,7},ymajorgrids=true]
\addplot[fill=IlliniOrange,draw=IlliniOrange] coordinates {(1,-2)(2,7)};
\draw[IlliniBlue] (axis cs:.6,0)--(axis cs:2.4,0);
\end{axis}
\end{tikzpicture}

\caption{High response and high incremental value can identify different customers.}\label{fig:M01_response_vs_value}
\end{figure}

With contribution and contact cost held fixed, the first account's high
baseline renewal probability reverses its ranking under incremental
value. A response model and a targeting policy therefore need separate
evaluation.

\hypertarget{M01-the-unit-of-observation-determines-the-calculation}{%
\section{The Unit of Observation Determines the
Calculation}\label{M01-the-unit-of-observation-determines-the-calculation}}

\noindent A transaction system often stores one row per product line, while the
marketing question concerns orders or customers. Suppose order O1
contains two product lines, O2 contains one, and O3 contains three.
There are six rows, three orders, and perhaps only two customers.
Counting rows as purchases would report six orders. Averaging customer
attributes over product lines would give more weight to customers with
larger baskets.

A key uniquely identifies the intended unit. At the line level, the key
might be the combination of order ID and line number. At the order
level, it is order ID. At the customer level, it is customer ID. A
useful data contract states the key, allowed duplicates, measurement
units, missing-value codes, and date conventions. It also states whether
a negative quantity represents a return, whether a price includes tax,
and whether canceled orders belong in the analysis.

For line \(\ell\) in order \(o\), gross line sales are quantity times
unit price, \(s_{o\ell}=q_{o\ell}p_{o\ell}\). Order sales are
\(S_o=\sum_\ell s_{o\ell}\), and customer-period sales are the sum
across that customer's eligible orders during the period. These are
identities once the population and variables have been defined. They do
not require a behavioral assumption. The substantive choices enter
through eligibility, returns, discounts, and the observation window.

Joining data can change the population or its weights. A many-to-one
join from order lines to a customer table should add customer attributes
without multiplying the original rows. If the customer table
accidentally contains two records for a customer, every matching order
line appears twice. A program can finish without an error while
overstating sales. Verify the relationship between keys before joining,
and compare row counts and totals afterward.

\begin{figure}[htbp]
\centering
% Original synthetic teaching figure; see corresponding module.
\begin{tikzpicture}
\node[box,text width=3.1cm] (o) at (0,0) {One order\\Customer C1\\Revenue: \$40};
\node[box,text width=3.4cm] (one) at (5,1.6) {Unique lookup key\\C1: one customer row};
\node[orangebox,text width=3.4cm] (two) at (5,-1.6) {Duplicate lookup key\\C1: two customer rows};
\node[box,text width=2.8cm] (good) at (10,1.6) {One joined row\\Total: \$40};
\node[orangebox,text width=2.8cm] (bad) at (10,-1.6) {Two joined rows\\Total: \$80};
\draw[edge] (o)--(one); \draw[accent] (o)--(two);
\draw[edge] (one)--(good); \draw[accent] (two)--(bad);
\node[align=center,text width=11cm] at (5,-3.2) {A valid-looking join can double revenue.\\Check lookup-key uniqueness before merging, then reconcile totals.};
\end{tikzpicture}

\caption{A duplicate customer key multiplies transactions in a join.}\label{fig:M01_join_multiplication}
\end{figure}

For a single \$40 order, a customer lookup with two matching records
produces two joined rows and an apparent \$80 total. This error arises
from the structure of the tables. Increasing the sophistication of the
subsequent model cannot repair it.

\hypertarget{M01-python-objects-and-reproducible-calculations}{%
\section{Python Objects and Reproducible
Calculations}\label{M01-python-objects-and-reproducible-calculations}}

\noindent Python stores individual values in objects such as integers,
floating-point numbers, strings, and booleans. A list stores an ordered
collection; a dictionary associates names with values. A function
packages a calculation with explicit inputs and an output. These
distinctions matter when an assistant supplies code: a column of
numeric-looking strings does not behave like a numeric column, and an
empty string does not automatically represent a missing value.

Suppose two orders contain three product lines. Multiplying quantity by
unit price gives line sales; summing within each order gives order
sales. Returning each calculation as a number allows its result to be
checked and reused.

\noindent\begin{minipage}{\linewidth}
\begin{Verbatim}[breaklines]
lines = [
    {"order": "O1", "quantity": 2, "unit_price": 10.0},
    {"order": "O1", "quantity": 1, "unit_price": 5.0},
    {"order": "O2", "quantity": 3, "unit_price": 4.0},
]

def line_sales(line):
    return line["quantity"] * line["unit_price"]

order_sales = {}
for line in lines:
    key = line["order"]
    order_sales[key] = order_sales.get(key, 0.0) + line_sales(line)

assert order_sales == {"O1": 25.0, "O2": 12.0}
assert sum(order_sales.values()) == 37.0
\end{Verbatim}
\end{minipage}

The total is \$37 across two orders, so average order value is \$18.50.
Dividing by the three product lines would instead give average line
sales, \$12.33. Neither denominator is inherently wrong; they describe
different objects. The mistake would be assigning one calculation the
other's interpretation.

A DataFrame provides a tabular interface for larger data. In pandas,
selecting columns, filtering rows, creating variables, grouping
observations, and joining tables are distinct operations. Use explicit
intermediate names when those operations change the unit of observation.
For example, \texttt{eligible\_lines}, \texttt{orders}, and
\texttt{customer\_months} communicate more than repeatedly overwriting
\texttt{df}. The
\href{https://pandas.pydata.org/docs/getting_started/intro_tutorials/}{pandas
tutorials} provide syntax references; the analytical specification
should determine which operations are needed.

Notebook execution order can hide dependencies. A cell may work because
an earlier interactive command created a variable that is absent from
the saved notebook. Before describing a notebook as reproducible,
restart the kernel and run all cells in order. Store file paths relative
to the project, record package versions, and distinguish source data
from derived outputs. A random seed helps reproduce a simulation or
split, but it does not guarantee identical results across every software
version or computing platform.

\hypertarget{M01-case-study-a-coupon-campaign-specification}{%
\section{Case Study: A Coupon Campaign
Specification}\label{M01-case-study-a-coupon-campaign-specification}}

\noindent A retailer has a monthly contact budget for 5,000 customers. The
proposed campaign offers a discount on the next purchase. The manager
initially asks for ``a model to find our best customers.'' That request
leaves the action, outcome, and objective unresolved. The
highest-spending customers may receive the offer without changing their
behavior; customers likely to respond may buy products with low margins.

A usable specification would identify the decision date, eligible
customer population, offer, evaluation horizon, and objective. One
possibility is to maximize incremental contribution during the next 30
days, subject to the contact limit. Eligibility might require marketing
consent and an active account. The analysis should define contribution
after discounts, returns, fulfillment, and contact costs. It should
identify which customer information existed before assignment and how
untreated outcomes will be measured.

The initial dataset might include customer-level histories and a
randomized pilot. Without such a pilot or another credible design, a
response model can help describe or predict behavior, but the
incremental contribution calculation remains unresolved. The next
operational step may therefore be a randomized test, rather than
immediate large-scale optimization.

A specification also needs acceptance criteria. The eligible-customer
table must have one row per customer, no post-assignment features, and
the intended date coverage. A campaign summary must reconcile assigned
counts, contact costs, purchases, and contribution. These checks are
small enough to state before asking an AI assistant to implement the
analysis. They protect the meaning of the result while allowing
flexibility in the code.

\hypertarget{M01-ai-assisted-coding-as-a-sequence-of-verifiable-steps}{%
\section{AI-Assisted Coding as a Sequence of Verifiable
Steps}\label{M01-ai-assisted-coding-as-a-sequence-of-verifiable-steps}}

\noindent An effective coding request names the inputs, output, definitions, and
checks. For the transaction example, ask for a function that accepts
product-line records, validates required fields, aggregates sales by
order, and returns both order totals and an exception report. Provide
the three-line example and its expected \$37 total. Ask the assistant to
explain what happens with a return, a missing order ID, or a duplicated
line.

Generate one analytical transformation at a time when the next step
depends on a definition. Inspect the proposed code before executing it,
especially commands that install software, delete files, upload data, or
invoke an external API. Read error messages as evidence about the code's
assumptions. A type error may reveal numeric values stored as strings; a
merge error may reveal an incorrect key. Repairing the error without
understanding its cause can preserve the underlying analytical mistake.

Verification should use both identities and meaningful counterexamples.
Order totals must sum to the line total for the same eligible rows. A
single known customer should produce a manually calculable summary. An
empty input should have an explicit result. A missing customer lookup
should not silently remove sales. These checks constrain what the
implementation is allowed to do without forcing one particular coding
style.

Keep a short record of consequential AI contributions: the task
requested, important changes to the plan or code, and the checks
performed. A saved conversation alone does not establish correctness.
The record should make it possible to explain why a transformation was
accepted.

\hypertarget{M01-interpreting-and-communicating-a-result}{%
\section{Interpreting and Communicating a
Result}\label{M01-interpreting-and-communicating-a-result}}

\noindent A useful recommendation names the decision, states the relevant estimate
with units and uncertainty, and explains how it changes the action.
``The model has high accuracy'' leaves the manager's decision
unresolved. ``At a \$3 contact cost and the stated contribution
assumptions, contact the customers with positive estimated incremental
contribution, subject to capacity'' connects the analysis to the
objective. The recommendation still needs evidence about how the
probabilities were estimated and whether the policy was evaluated on new
customers or periods.

Results should distinguish observation from assumption. A transaction
total is calculated from recorded eligible rows. A renewal probability
is an estimate. A contribution margin may be supplied by finance. An
untreated outcome for an offered customer is counterfactual. Listing
these roles prevents a spreadsheet or notebook from giving every input
the appearance of equal certainty.

If one customer places four orders and another places one, an
order-weighted mean gives the first customer four times the weight. If
both customers should count equally, calculate each customer's mean
first. If an AI-generated program reports six purchases from six product
lines, inspect whether purchase was defined as a line or an order. If an
offer has a high response rate, ask for the response rate under no offer
before interpreting the response as an effect.

\hypertarget{M01-readings}{%
\section{Readings}\label{M01-readings}}

\noindent The
\href{https://github.com/python-marketing-research/python-marketing-research-1ed}{Python
marketing research companion repository} provides examples tied to
marketing applications. Its notebooks are reference material and require
inspection of dependencies and data paths before adaptation. The
\href{https://pandas.pydata.org/docs/getting_started/intro_tutorials/}{pandas
tutorials} support tabular implementation.

\hypertarget{M02-m2.-customer-data-measurement-and-visualization}{%
\chapter{Customer Data, Measurement, and
Visualization}\label{M02-m2.-customer-data-measurement-and-visualization}}

\noindent A grocery retailer sees larger online baskets and asks whether moving
customers to its app will raise sales. The raw comparison combines the
shopping channel with who uses it, why they shop, and which trips are
observed. A household may order online for a planned weekly restock and
visit a store for a forgotten ingredient. Measurement and visualization
determine whether the analysis reveals that distinction or conceals it.

\hypertarget{M02-define-the-population-before-the-metric}{%
\section{Define the Population Before the
Metric}\label{M02-define-the-population-before-the-metric}}

\noindent A customer dataset is a record of an observation process. Loyalty
accounts exclude unidentified shoppers. Online logs record sessions that
may contain repeated visits by the same person. Orders exclude browsing
without purchase. App events reflect tracking rules, consent, device
settings, and instrumentation changes. Before calculating a metric,
identify the target population and the mechanism by which its members
enter the table.

Suppose monthly revenue per customer is intended to describe all
eligible accounts. Dividing total revenue by purchasing accounts
produces revenue per purchasing account. If inactivity changes, that
measure can increase while revenue per eligible account falls. The
denominator should therefore be written into the metric name or its
definition. For customer acquisition cost, specify whether the
denominator is new registered accounts, first-time purchasers, or
customers with positive contribution over a defined horizon.

A conversion rate requires an eligible exposure count. If \(b\) eligible
sessions produce at least one purchase among \(n\) eligible sessions,
session conversion is \(b/n\). Customer conversion instead uses
customers in both numerator and denominator. A customer with ten
sessions should not receive ten times the weight in a customer-level
rate. When the same customer contributes repeated observations,
statistical uncertainty must also account for that dependence.

Returns require a substantive rule. Gross sales, net sales after
returns, and contribution after variable costs answer different
questions. A comparison that deducts online returns but observes only
gross offline sales mixes measurement rules. Currency, tax, fulfillment
fees, and discounts should be aligned across channels before
interpreting differences.

\hypertarget{M02-build-a-customer-table-without-losing-customers}{%
\section{Build a Customer Table Without Losing
Customers}\label{M02-build-a-customer-table-without-losing-customers}}

\noindent Begin with a customer roster that defines eligibility. Aggregate
eligible transactions to customer-period totals, then left-join those
totals to the roster. A customer with no transactions can receive zero
sales if the transaction system fully captures activity during the
observation period. Zero should not replace a value that is missing
because the customer was unobserved or the data feed failed.

A join should have a documented relationship: one-to-one, many-to-one,
or one-to-many. Validate key uniqueness in the lookup table and inspect
unmatched records. A match rate of 99\% can still conceal a material
problem if the unmatched 1\% contains the highest-value customers or one
entire channel. Report both counts and relevant value totals for
unmatched observations.

Dates need similar attention. Convert timestamps to a consistent
timezone before assigning day or week. Determine whether the date
represents ordering, shipment, payment, or return. A promotion announced
late in a period may affect orders immediately but recognized revenue
later. Aligning metrics on an inappropriate date can create an apparent
lag that comes from accounting rather than customer behavior.

Outliers should be investigated before removal. A \$10,000 order could
be a business customer, a data-entry error, or a legitimate event. A
rule that trims all orders above a percentile changes the estimand to a
restricted population. State that restriction and compare conclusions
with and without the rule when it can affect the decision. Robust
summaries such as medians and quantiles often reveal more than one mean
with an undocumented exclusion.

\hypertarget{M02-case-study-a-reversal-caused-by-customer-mix}{%
\section{Case Study: A Reversal Caused by Customer
Mix}\label{M02-case-study-a-reversal-caused-by-customer-mix}}

\noindent A marketing manager compares campaigns A and B using conversion. Both
campaigns reach new and returning customers, but their audience
composition differs. Assume the following customer counts and outcomes,
with each customer belonging to one segment and appearing once.

\begin{table}[H]
\centering\small
\begin{tabular}{@{}
  >{\raggedright\arraybackslash}p{(\columnwidth - 10\tabcolsep) * \real{0.13}}
  >{\raggedleft\arraybackslash}p{(\columnwidth - 10\tabcolsep) * \real{0.17}}
  >{\raggedleft\arraybackslash}p{(\columnwidth - 10\tabcolsep) * \real{0.17}}
  >{\raggedleft\arraybackslash}p{(\columnwidth - 10\tabcolsep) * \real{0.17}}
  >{\raggedleft\arraybackslash}p{(\columnwidth - 10\tabcolsep) * \real{0.17}}
  >{\raggedleft\arraybackslash}p{(\columnwidth - 10\tabcolsep) * \real{0.19}}@{}}
\toprule
Campaign & New Customers & New Conversions & Returning Customers &
Returning Conversions & Overall Conversion \\
\midrule

A & 100 & 9 & 100 & 80 & 89/200 = 44.5\% \\
B & 200 & 20 & 100 & 81 & 101/300 = 33.7\% \\
\bottomrule
\end{tabular}
\end{table}

Campaign B has higher conversion among new customers, 10\% versus 9\%,
and among returning customers, 81\% versus 80\%. Campaign A nevertheless
has the higher aggregate rate because half its audience consists of
returning customers, compared with one-third for B. Returning customers
convert much more frequently in both campaigns.

The aggregate rate is a weighted average of segment rates:

\[
\bar p_j=\sum_g w_{jg}p_{jg},\qquad \sum_g w_{jg}=1.
\]

Here \(j\) indexes campaigns, \(g\) indexes customer segments,
\(p_{jg}\) is the conversion fraction within a segment, and \(w_{jg}\)
is that segment's share of the campaign audience. The weights are part
of the comparison. Holding them at 50\% new and 50\% returning gives
standardized rates of 44.5\% for A and 45.5\% for B. This adjustment
describes performance for the chosen common mix. It establishes a causal
campaign effect only if the design supports that interpretation.

\begin{figure}[htbp]
\centering
% Original synthetic teaching figure; see corresponding module.
\begin{tikzpicture}
\begin{axis}[ybar,bar width=15pt,width=13cm,height=7cm,ymin=0,ymax=100,xmin=.4,xmax=4.6,xtick={1,2,3,4},xticklabels={New,Returning,Observed mix,Common mix},ylabel={Conversion rate (percent)},ytick={0,20,40,60,80,100},ymajorgrids=true,legend style={at={(.5,1.13)},anchor=south,legend columns=2},nodes near coords,nodes near coords style={font=\scriptsize},every node near coord/.append style={/pgf/number format/fixed,/pgf/number format/precision=1}]
\addplot[fill=IlliniBlue,draw=IlliniBlue] coordinates {(1,9)(2,80)(3,44.5)(4,44.5)};
\addplot[fill=IlliniOrange,draw=IlliniOrange] coordinates {(1,10)(2,81)(3,33.6667)(4,45.5)};
\legend{Campaign A,Campaign B}
\end{axis}
\end{tikzpicture}

\caption{Audience composition reverses the aggregate campaign ranking.}\label{fig:M02_mix_reversal}
\end{figure}

B's small advantages within both segments coexist with A's aggregate
advantage. A manager deciding where to send the next campaign needs to
specify the future audience mix and whether customers were assigned
comparably. An aggregate bar chart alone cannot answer that question.

\hypertarget{M02-cohorts-retention-and-customer-age}{%
\section{Cohorts, Retention, and Customer
Age}\label{M02-cohorts-retention-and-customer-age}}

\noindent A cohort groups customers by a common starting event, such as first
purchase or subscription activation. Retention at customer age \(k\) is
the fraction of the original eligible cohort meeting a stated activity
criterion at that age. A cohort table separates calendar time from
customer age: an October signup observed one month later is not
comparable with an April signup observed seven months later unless the
question explicitly concerns those different ages.

For cohort \(c\), define \(N_c\) as the number of eligible customers at
entry and \(A_{c,k}\) as the number active at age \(k\). The descriptive
retention rate is \(R_{c,k}=A_{c,k}/N_c\). A later cohort may not yet
have reached age \(k\); that cell is unobserved, not zero. Omitting
inactive customers from the denominator would turn retention into a
different conditional measure.

\begin{table}[H]
\centering\small
\begin{tabular}{@{}
  >{\raggedright\arraybackslash}p{(\columnwidth - 8\tabcolsep) * \real{0.15}}
  >{\raggedleft\arraybackslash}p{(\columnwidth - 8\tabcolsep) * \real{0.21}}
  >{\raggedleft\arraybackslash}p{(\columnwidth - 8\tabcolsep) * \real{0.21}}
  >{\raggedleft\arraybackslash}p{(\columnwidth - 8\tabcolsep) * \real{0.21}}
  >{\raggedleft\arraybackslash}p{(\columnwidth - 8\tabcolsep) * \real{0.22}}@{}}
\toprule
Signup Cohort & Entry Customers & Active at Month 1 & Active at Month 2
& Active at Month 3 \\
\midrule

January & 100 & 60 & 45 & 40 \\
February & 100 & 70 & 50 & Not yet observed \\
March & 100 & 80 & Not yet observed & Not yet observed \\
\bottomrule
\end{tabular}
\end{table}

Month-one retention rises from 60\% for the January cohort to 80\% for
the March cohort. March customers have not yet reached the later
observation windows, so their long-run retention remains unobserved.
Comparing March's month one with January's month three would combine
cohort quality with customer age. The right comparison depends on
whether the decision concerns onboarding, long-term value, or a
calendar-time change affecting all existing customers.

\begin{figure}[htbp]
\centering
% Original synthetic teaching figure; see corresponding module.
\begin{tikzpicture}[x=2.25cm,y=1.05cm]
\foreach \c/\lab in {0/At Entry,1/Month 1,2/Month 2,3/Month 3} {\node[font=\small\bfseries] at (\c,1) {\lab};}
\foreach \r/\lab in {0/January,-1/February,-2/March} {\node[anchor=east] at (-.6,\r) {\lab};}
\foreach \x/\y/\v in {0/0/100,1/0/60,2/0/45,3/0/40,0/-1/100,1/-1/70,2/-1/50,0/-2/100,1/-2/80} {\draw[IlliniBlue] (\x-.47,\y-.4) rectangle (\x+.47,\y+.4);\node at (\x,\y) {\v\%};}
\foreach \x/\y in {3/-1,2/-2,3/-2} {\fill[LightStorm] (\x-.47,\y-.4) rectangle (\x+.47,\y+.4);\node[align=center,font=\scriptsize] at (\x,\y) {Not yet\\observed};}
\node[align=center,text width=11cm] at (1.3,-3.2) {Each cohort begins with 100 customers.\\Only observed customer ages enter the retention denominator.};
\end{tikzpicture}

\caption{A cohort’s unobserved future is different from zero retention.}\label{fig:M02_cohort_triangle}
\end{figure}

Unobserved future periods require forecasts from an explicit model.
Coding those periods as zero would treat unavailable outcomes as
observed failures and bias retention downward.

\hypertarget{M02-visual-encodings-should-match-the-question}{%
\section{Visual Encodings Should Match the
Question}\label{M02-visual-encodings-should-match-the-question}}

\noindent Use position on a common scale when precise comparison matters. A dot
plot can compare segment conversion rates more clearly than area-sized
bubbles. A line chart is appropriate when the ordering of time is
meaningful. A histogram or empirical cumulative distribution reveals
skewness that a mean hides. A scatterplot supports examination of a
relationship, with transparency or aggregation when observations
overlap.

The scale changes the interpretation. A logarithmic sales axis
emphasizes proportional differences; a linear axis emphasizes dollar
differences. A truncated bar-chart axis exaggerates differences because
bar length encodes magnitude from its baseline. A carefully labeled
narrow axis can be useful for a line or dot plot when the scientific
question concerns small differences. The choice should preserve the
intended comparison and disclose the scale.

Uncertainty should be shown for an estimated comparison when it affects
a decision. For an independent Bernoulli sample, a common large-sample
standard error for a proportion is \(\sqrt{\hat p(1-\hat p)/n}\). Near
zero or one, or with small samples, use an interval appropriate to
binomial data. With repeated observations per customer, treating every
session as independent generally understates uncertainty. Resample or
cluster at the unit that captures the relevant dependence.

Color can encode categories or highlight the comparison that matters.
Use line styles, markers, and labels alongside color, and keep small
text dark enough for contrast. A figure should remain interpretable when
printed in grayscale or described to a student who cannot see it.
Explain the axes, units, sample, and principal comparison in its caption
or accompanying text.

\hypertarget{M02-case-study-online-grocery-baskets-and-repeated-purchases}{%
\section{Case Study: Online Grocery Baskets and Repeated
Purchases}\label{M02-case-study-online-grocery-baskets-and-repeated-purchases}}

\noindent Chintala, Liaukonytė, and Yang's
\href{https://doi.org/10.1287/mksc.2022.0292}{grocery-shopping study}
compares online and offline baskets using household shopping records. It
studies basket composition and repeated purchasing, with attention to
differences in shopping-trip types. Larger baskets need not contain a
wider range of categories, and a household's choice of channel can be
related to the purpose of its trip.

Suppose one household buys the sets \{milk, bread, apples, yogurt\} and
\{milk, bread, apples, coffee\} on successive trips. Each basket
contains four distinct items. Their Jaccard similarity is

\[
J(A,B)=\frac{|A\cap B|}{|A\cup B|}=\frac{3}{5}=0.60.
\]

The numerator counts shared distinct items and the denominator counts
all distinct items appearing in either trip. A different pair of
four-item baskets with one shared item has similarity
\(1/7\approx0.143\). Basket size is identical in the two examples, while
repeated purchasing differs substantially. Whether quantities or
spending weights should enter the measure depends on the decision;
set-based Jaccard treats one unit and ten units of an item alike.

A defensible descriptive study would define eligible grocery trips,
compare channels within households when appropriate, and examine basket
size, distinct categories, repeat-item overlap, and spending. It would
inspect whether findings are driven by large restocking trips, small
convenience trips, or particular retailers. Channel labels should not
silently determine the definition of a ``typical'' trip when the
objective is a comparable classification. Clustering can classify trips
using their observed composition.

A retail recommendation might distinguish discovery from replenishment.
High repetition can support convenient reordering while making it harder
for unfamiliar products to receive attention. A proposed interface
change should therefore specify both the desired outcome and the
possible tradeoff. Its effect requires evaluation under an appropriate
design, even when the descriptive pattern provides a strong reason to
test it.

\hypertarget{M02-from-evidence-to-a-management-presentation}{%
\section{From Evidence to a Management
Presentation}\label{M02-from-evidence-to-a-management-presentation}}

\noindent The Pyramid Principle organizes a recommendation around a
decision-relevant answer, followed by a small number of supporting
reasons and the evidence for each reason. This structure works when the
answer's strength matches the analysis. A descriptive study might
recommend a test or a closer diagnostic comparison, rather than
immediate deployment.

For the campaign example, a clear opening is: ``Compare the campaigns at
a common customer mix before reallocating the budget.'' The supporting
reasons are that returning customers convert more often, the campaigns
reached different mixes, and the within-segment ranking differs from the
aggregate ranking. The supporting evidence should quantify each reason
and connect it to the allocation decision.

An effective chart title can state the comparison, such as ``Campaign
A's aggregate advantage reflects its returning-customer share.'' The
caption should define the conversion denominator and identify the data
as synthetic.

\hypertarget{M02-ai-assisted-workflow-and-practice}{%
\section{AI-Assisted Workflow and
Practice}\label{M02-ai-assisted-workflow-and-practice}}

\noindent Ask an assistant to construct a customer-period table from a specified
roster and transaction table, preserving customers with no observed
purchases. Require a join audit, a missing-data report, and a
reconciliation of total sales. Next ask it to calculate the campaign
table above and standardized rates under a 50/50 audience mix. Verify
those small numbers before applying the same functions to a larger file.

For a chart request, specify the question, observation unit, axis units,
ordering, and uncertainty representation. Ask the assistant to describe
what visual pattern would change the recommendation. Compare its
description with the plotted data. If it labels an observational channel
difference as an effect, revise the interpretation and identify the
design that would be needed to support the stronger claim.

For practice, change campaign B's audience to 100 new and 100 returning
customers while keeping its segment conversion rates fixed. Its
aggregate rate becomes 45.5\%, and the reversal disappears. Next
calculate month-two retention for the January and February cohorts: 45\%
and 50\%. March has no observed month-two rate. Finally, compare a
four-item basket with itself: Jaccard similarity is one, regardless of
whether the quantities have changed. Each exercise tests a definition
that should remain explicit in an automated analysis.

\hypertarget{M02-readings}{%
\section{Readings}\label{M02-readings}}

\noindent \href{https://clauswilke.com/dataviz/}{Wilke's visualization text}
offers a visual reference for distributions, comparisons, and
uncertainty. The
\href{https://pandas.pydata.org/docs/user_guide/groupby.html}{pandas
grouping documentation} supports implementation.

\hypertarget{M03-m3.-regression-interpretation-and-marketing-decisions}{%
\chapter{Regression, Interpretation, and Marketing
Decisions}\label{M03-m3.-regression-interpretation-and-marketing-decisions}}

\noindent A retailer observes that stores with larger advertising budgets sell
more. The manager wants to know whether another dollar of advertising
would be profitable. The observed relationship combines the response to
advertising with differences in store size, local demand, seasonality,
and the firm's reasons for choosing the budget. Regression provides a
way to summarize conditional relationships. Its value depends on which
variation identifies the comparison and whether that comparison answers
the manager's question.

\hypertarget{M03-a-conditional-mean-and-a-regression-line}{%
\section{A Conditional Mean and a Regression
Line}\label{M03-a-conditional-mean-and-a-regression-line}}

\noindent Let \(Y_i\) be weekly store sales measured in thousands of dollars and
\(X_i\) be weekly advertising expenditure in thousands of dollars. A
simple linear model writes

\[
Y_i=\beta_0+\beta_1X_i+\varepsilon_i.
\]

The intercept \(\beta_0\) is the fitted sales level at zero advertising;
its practical interpretation is weak when zero lies far outside the
observed range. The slope \(\beta_1\) is the change in fitted sales, in
thousands of dollars, associated with an additional thousand dollars of
advertising. The residual captures the difference between observed sales
and the fitted line. Linearity is an approximation to the relationship
over the relevant range, not a claim that every customer's behavior is
linear.

Ordinary least squares chooses coefficients to minimize the sum of
squared residuals. In a simple regression with an intercept, its slope
is

\[
\hat\beta_1=\frac{\sum_i(X_i-\bar X)(Y_i-\bar Y)}{\sum_i(X_i-\bar X)^2}.
\]

The numerator measures whether high-advertising observations tend to
have high sales. The denominator measures how much advertising varies.
When advertising hardly varies, the data provide little information
about a slope, even if the sample contains many rows. A large sample of
nearly identical budgets does not substitute for informative variation.

For three observations \((X,Y)=(1,3),(2,5),(3,7)\), the fitted line is
\(\hat Y=1+2X\). At \(X=2.5\), it predicts sales of \$6,000. If the
sales contribution margin is 30\%, an additional \$1,000 of advertising
is associated with \$600 of contribution from \$2,000 of extra sales.
Treating that association as a return to intervention requires an
identification argument. Even under a causal interpretation, the
incremental contribution is less than the additional advertising expense
in this numerical example.

\hypertarget{M03-multiple-regression-and-the-comparison-it-uses}{%
\section{Multiple Regression and the Comparison It
Uses}\label{M03-multiple-regression-and-the-comparison-it-uses}}

\noindent A multiple regression can include observed store characteristics:

\[
Y_i=\beta_0+\beta_1X_i+\gamma Z_i+\varepsilon_i,
\]

where \(Z_i\) might be floor area measured in thousands of square feet.
The advertising coefficient summarizes variation in sales associated
with advertising after removing the linear association with floor area.
A useful way to understand this is to regress advertising on the other
controls, regress sales on the same controls, and regress the resulting
sales residuals on the advertising residuals. The Frisch-Waugh-Lovell
result gives the same advertising coefficient as the full regression
when the residualization includes the same intercept and controls.

``Controlling for'' a variable therefore changes the comparison. It does
not automatically remove all confounding. If managers observe local
demand information that the analyst lacks, their advertising choices may
remain related to unobserved sales determinants. Controls measured after
the intervention can also remove part of the effect or introduce bias.
Select variables using the decision timing and a substantive account of
how the data were generated.

When an omitted variable affects both the predictor and the outcome, the
simple slope can be misleading. In the linear population model above,
suppose the regression disturbance is mean-independent of \(X\) and
\(Z\). Omitting \(Z\) yields a simple-regression slope equal to

\[
\beta_1+\gamma\frac{\operatorname{Cov}(X,Z)}{\operatorname{Var}(X)}.
\]

The extra term has a substantive interpretation. If larger stores sell
more, so \(\gamma>0\), and larger stores spend more on advertising, so
the covariance is positive, the simple advertising slope exceeds the
conditional slope. The equation is a model-based decomposition; its
interpretation depends on the stated linear specification and exogeneity
condition.

\hypertarget{M03-case-study-attractive-markets-and-apparent-competitive-benefits}{%
\section{Case Study: Attractive Markets and Apparent Competitive
Benefits}\label{M03-case-study-attractive-markets-and-apparent-competitive-benefits}}

\noindent \href{https://doi.org/10.3982/QE478}{Igami and Yang (2016)} study entry
and competition among hamburger chains in Canada. Their setting
motivates a general identification problem: attractive markets can
support both more competitors and better outcomes for a focal firm. A
positive relationship between rival presence and performance can
therefore coexist with a negative competitive effect.

Within low-demand markets, suppose stores facing one, two, and three
rivals have sales of 9, 8, and 7 units. Within high-demand markets,
stores facing four, five, and six rivals have sales of 16, 15, and 14
units. Within each demand group, one additional rival is associated with
one unit less sales. Across the six observations, the pooled slope is
approximately \(+1.57\): high-demand markets have more rivals and higher
sales.

\begin{figure}[htbp]
\centering
% Original synthetic teaching figure; see corresponding module.
\begin{tikzpicture}
\begin{axis}[xmin=.5,xmax=6.5,ymin=5,ymax=18,xlabel={Number of rivals},ylabel={Sales (illustrative units)},legend style={at={(.5,1.08)},anchor=south,legend columns=2},xtick={1,...,6},ymajorgrids=true]
\addplot[IlliniBlue,mark=*,very thick] coordinates {(1,9)(2,8)(3,7)};
\addplot[IlliniOrange,mark=square*,very thick] coordinates {(4,16)(5,15)(6,14)};
\addplot[Storm,dashed,domain=1:6] {6+11*x/7};
\legend{Low demand,High demand,Pooled fit}
\node[anchor=west,font=\small] at (axis cs:1.1,6.3) {Within-group slope: $-1$};
\node[anchor=west,font=\small] at (axis cs:3.4,17) {Pooled slope: $+1.57$};
\end{axis}
\end{tikzpicture}

\caption{A positive pooled relationship can conceal a negative within-market-type relationship.}\label{fig:M03_omitted_demand}
\end{figure}

The within-group slope is \(-1\), while baseline demand differs between
groups. The positive pooled slope reflects that demand composition. In
an actual retail application, market demand may be unobserved, entry may
be forward-looking, and firms may respond strategically.

A manager considering entry should ask which outcomes would change if an
additional competitor entered the same market. Comparing different
markets does not automatically provide that counterfactual. Market fixed
effects, measured demand controls, an instrument, or a credible event
design may improve the comparison, but each uses different variation and
requires its own conditions. A highly predictive regression can remain
unsuitable for interpreting the coefficient as competition's causal
effect.

\hypertarget{M03-transformations-and-economically-meaningful-units}{%
\section{Transformations and Economically Meaningful
Units}\label{M03-transformations-and-economically-meaningful-units}}

\noindent A level-level specification gives an additive interpretation: one more
unit of \(X\) changes fitted \(Y\) by \(\beta_1\). A log-level
specification, \(\log Y=\beta_0+\beta_1X+\varepsilon\), describes
proportional differences in a positive outcome. A one-unit change in
\(X\) changes the fitted log outcome by \(\beta_1\); the corresponding
multiplicative change in the log-scale prediction is \(e^{\beta_1}\), or
\(100(e^{\beta_1}-1)\) percent. The approximation \(100\beta_1\) percent
is useful when the coefficient is small.

A log-log demand model, \(\log Q=\alpha+\beta\log P+\varepsilon\), gives
\(\beta\) an elasticity interpretation over the range where the
specification applies. If \(\beta=-1.5\), a 1\% higher price is
associated with approximately 1.5\% lower quantity. Sales revenue is
\(PQ\), so its local price elasticity is \(1+\beta=-0.5\). Profit also
depends on unit cost. A price that raises revenue can still lower
profit, and observed prices can be endogenous because firms set them in
response to demand.

Zero and negative outcomes need deliberate treatment. Taking a logarithm
of zero is undefined. Adding one changes the scale and interpretation,
especially when values are small or the unit of measurement changes. A
two-part model can separate purchase incidence from positive
expenditure. A suitable generalized linear model may model the
conditional mean directly. The choice should reflect the outcome and
intended estimand, with diagnostics for the selected specification.

Back-transforming a predicted log outcome also needs care.
Exponentiating the predicted conditional mean of \(\log Y\) generally
does not produce the conditional mean of \(Y\), because exponentiation
is nonlinear. Retransformation methods or a model for the mean on the
original scale may be needed when total sales is the target. Evaluate
predictions in the units used for the managerial decision.

\hypertarget{M03-interactions-describe-how-a-relationship-changes}{%
\section{Interactions Describe How a Relationship
Changes}\label{M03-interactions-describe-how-a-relationship-changes}}

\noindent Suppose an analyst expects a promotional message to work differently for
customers with long tenure. A linear interaction specification is

\[
Y_i=\beta_0+\beta_1D_i+\beta_2T_i+\beta_3D_iT_i+\varepsilon_i,
\]

where \(D_i\) is a message indicator and \(T_i\) is tenure in years. The
fitted message difference at tenure \(t\) is \(\beta_1+\beta_3t\). The
coefficient \(\beta_1\) is the difference at zero tenure, so centering
tenure at a meaningful value can make that coefficient easier to
interpret. The interaction coefficient tells how the fitted difference
changes per year of tenure.

With \(\beta_1=4\) dollars and \(\beta_3=-1\) dollar per year, the
fitted difference is \$4 at tenure zero, \$2 at two years, and zero at
four years. An average coefficient can conceal that variation.
Predictions beyond the observed tenure range should be treated as
extrapolation. In a nonlinear probability model, an interaction
coefficient on the model's index does not directly equal the interaction
in probability; compute and plot the predicted probabilities or marginal
contrasts.

\begin{figure}[htbp]
\centering
% Original synthetic teaching figure; see corresponding module.
\begin{tikzpicture}
\begin{axis}[xmin=0,xmax=5,ymin=-1.6,ymax=5,xlabel={Customer tenure (years)},ylabel={Fitted message difference (dollars)},xtick={0,...,5},ytick={-1,0,1,2,3,4},ymajorgrids=true]
\addplot[IlliniOrange,very thick,domain=0:5] {4-x};
\addplot[IlliniBlue,dashed] coordinates {(0,0)(5,0)};
\addplot[IlliniBlue,only marks,mark=*] coordinates {(0,4)(2,2)(4,0)};
\node[anchor=south west] at (axis cs:.1,4) {$\Delta(t)=4-t$};
\node[anchor=south west,align=left] at (axis cs:3.3,.6) {Zero contrast\\at four years};
\draw[edge] (axis cs:4,.6)--(axis cs:4,.1);
\end{axis}
\end{tikzpicture}

\caption{An interaction changes the relevant comparison across customer tenure.}\label{fig:M03_interaction}
\end{figure}

For tenure between zero and five years, the assumed contrast \(4-t\)
crosses zero at four years. In an estimated application, uncertainty
around the contrast and support at each tenure would determine whether
such a crossing is informative.

\hypertarget{M03-uncertainty-dependence-and-practical-importance}{%
\section{Uncertainty, Dependence, and Practical
Importance}\label{M03-uncertainty-dependence-and-practical-importance}}

\noindent A standard error describes sampling uncertainty under an estimator and
its assumptions. A confidence interval provides a range generated by a
procedure with a stated repeated-sampling coverage property. A p-value
assesses compatibility with a particular null under its reference model.
None of these quantities measures the managerial value of an effect by
itself. A precisely estimated increase of one cent may be unimportant; a
wide interval around a large potential benefit may justify another
experiment.

For an approximately normal coefficient estimate, a rough 95\% interval
is the estimate plus or minus 1.96 standard errors, with appropriate
adjustments when a finite-sample reference distribution is needed. If a
sales slope is 2.0 with a standard error of 0.5, that interval is
approximately {[}1.02, 2.98{]} in the slope's units. Report the units
and the comparison alongside the interval.

Heteroskedasticity means that residual variance differs across
observations. Heteroskedasticity-robust standard errors can address that
feature for inference under suitable sampling conditions; they do not
repair omitted-variable bias. Repeated store-week observations may share
persistent store shocks. Clustered standard errors allow dependence
within a specified cluster under their asymptotic conditions. The
relevant number of independent clusters can be much smaller than the
number of rows. With few clusters, conventional approximations may be
unreliable and design-specific methods are needed.

\hypertarget{M03-diagnostics-and-panel-data}{%
\section{Diagnostics and Panel
Data}\label{M03-diagnostics-and-panel-data}}

\noindent Inspect residuals against fitted values and important predictors for
nonlinear patterns and changing variance. Investigate influential
observations, errors in units, and unusual leverage. Multicollinearity
can make individual coefficients unstable even when predictions are
reasonably stable, because several predictors explain nearly the same
variation. Removing a variable solely to obtain a smaller p-value
changes the question and can introduce omitted-variable bias.

Store fixed effects absorb time-invariant differences between stores.
Calendar fixed effects absorb common period-specific shocks. In a
store-week panel, a model with both uses changes within stores relative
to common time movements. It still requires a valid account of
time-varying confounders. If a store raises advertising precisely when
an unobserved local event increases demand, the advertising coefficient
can remain biased.

A fitted model should be evaluated on the task for which it is used. For
prediction, compare held-out errors with a useful baseline. For a
descriptive coefficient, inspect support and sensitivity to
substantively justified specifications. For a causal effect, assess the
design and identifying conditions. A large \(R^2\) indicates in-sample
fit relative to an intercept-only benchmark; it does not establish a
credible intervention comparison.

\hypertarget{M03-ai-assisted-implementation-and-practice}{%
\section{AI-Assisted Implementation and
Practice}\label{M03-ai-assisted-implementation-and-practice}}

\noindent Ask an assistant to fit the pooled and demand-group-adjusted regressions
in the six-observation example. Require a table with variable
definitions, units, coefficients, and the exact sample. Verify that the
adjusted rival coefficient is \(-1\) and the pooled coefficient is
\(11/7\approx1.57\). Ask it to explain why adding demand-group
indicators changes the comparison. A numerical match without that
explanation is an incomplete interpretation.

For a larger dataset, specify the formula and desired covariance
estimator before requesting code. Ask the assistant to list dropped rows
and explain how categorical variables were encoded. Have it calculate a
predicted contrast at two values of a moderator and compare the result
with the interaction formula. Treat diagnostic recommendations as
hypotheses to inspect, rather than automatic permission to remove
observations or alter the model.

For practice, consider an advertising coefficient measured per \$1,000
of spending. Re-express the predictor in dollars and explain how the
numerical coefficient changes while fitted sales remain the same. Next
explain why a robust standard error cannot eliminate the bias in the
market-demand example. Finally, distinguish a confidence interval for
the conditional mean from a prediction interval for a new store's
realized sales: the latter also reflects individual outcome variation.

\hypertarget{M03-readings}{%
\section{Readings}\label{M03-readings}}

\noindent \href{https://www.statlearning.com/}{An Introduction to Statistical
Learning} provides a foundation for regression, model assessment, and
nonlinear extensions. The
\href{https://www.statsmodels.org/stable/regression.html}{statsmodels
regression documentation} supports implementation.

\hypertarget{M04-m4.-product-preferences-customer-segments-and-lifetime-value}{%
\chapter{Product Preferences, Customer Segments, and Lifetime
Value}\label{M04-m4.-product-preferences-customer-segments-and-lifetime-value}}

\noindent A subscription retailer must choose a product bundle, decide whom to
serve, and determine how much to spend acquiring or retaining customers.
These decisions are connected. A feature that attracts one segment may
be costly to deliver, and a segment with high initial spending may have
low retention. Conjoint analysis estimates tradeoffs among product
attributes, clustering summarizes patterns of similarity, and customer
lifetime value translates future behavior into expected contribution.

\hypertarget{M04-product-design-requires-tradeoffs}{%
\section{Product Design Requires
Tradeoffs}\label{M04-product-design-requires-tradeoffs}}

\noindent Asking consumers whether they like fast delivery often produces a
favorable response, but it does not reveal how much they will pay for it
or which other feature they will give up. Conjoint tasks present product
profiles whose attributes vary. Respondents may rate profiles, rank
them, or choose among alternatives. Choice-based tasks resemble an
assortment decision, while ratings can provide more observations per
respondent under a different response process.

Attributes should represent decisions the firm can implement and
descriptions respondents can understand. Levels should span a plausible
range without creating impossible combinations. If every premium-quality
product also has the highest price, the design provides little
independent information about the quality-price tradeoff. Randomization
and a suitable experimental design help separate attribute effects, but
realistic constraints may limit which combinations can be shown.

Include the decision's relevant alternatives. A forced choice between
two products estimates preference conditional on choosing one of them.
An outside option such as ``buy neither'' helps represent category
participation. A simulation that omits the outside option can overstate
category demand because every respondent must buy something. Survey
design, task burden, and hypothetical behavior also affect how well
stated choices transfer to market behavior.

\hypertarget{M04-utility-choice-probabilities-and-willingness-to-pay}{%
\section{Utility, Choice Probabilities, and Willingness to
Pay}\label{M04-utility-choice-probabilities-and-willingness-to-pay}}

\noindent A simple model assigns respondent \(i\) a systematic utility for
alternative \(j\):

\[
V_{ij}=\beta_q q_j+\beta_f f_j+\beta_p p_j,
\]

where \(q_j\) indicates premium quality, \(f_j\) indicates a delivery
feature, and \(p_j\) is price in dollars. Utility has an arbitrary
scale. Coefficients describe tradeoffs on that scale, and price is
expected to enter negatively in a conventional demand interpretation.
Alternative-specific constants can capture average preference for named
options relative to a reference alternative.

Under a multinomial-logit random-utility specification, the probability
of choosing \(j\) from the available set \(\mathcal J_i\) is

\[
P_{ij}=\frac{\exp(V_{ij})}{\sum_{k\in\mathcal J_i}\exp(V_{ik})}.
\]

Multinomial logit imposes proportional substitution across alternatives.
Whether this pattern fits the product category requires empirical
assessment. Repeated choices by one respondent are dependent; estimation
and validation should respect that respondent structure. Heterogeneous
preferences may be represented with interactions, latent classes, or
random coefficients when the data support the additional complexity.

For a binary feature whose coefficient is \(\beta_f\), a local
willingness-to-pay tradeoff under linear price utility is
\(-\beta_f/\beta_p\). If \(\beta_f=1.2\) and \(\beta_p=-0.1\) per
dollar, the model equates the feature with \$12 of price. This ratio is
conditional on the specification and utility scale cancellation. A price
coefficient close to zero makes the ratio unstable, and preferences
outside the experimental price range are extrapolations.

Attribute importance is often calculated from the range of an
attribute's estimated part-worths relative to the sum of ranges. The
result depends on the levels included in the study. Expanding the price
range can increase measured price importance even when underlying
preferences are unchanged. Report the design range when discussing
importance.

\hypertarget{M04-case-study-a-new-bundle-can-cannibalize-existing-sales}{%
\section{Case Study: A New Bundle Can Cannibalize Existing
Sales}\label{M04-case-study-a-new-bundle-can-cannibalize-existing-sales}}

\noindent Consider two subscription bundles and an outside option. Bundle A has
premium quality and costs \$10; bundle B has standard quality and costs
\$6. Let \(\beta_q=1.2\), \(\beta_p=-0.1\), and normalize outside-option
utility to zero. The systematic utilities are 0.2 for A and \(-0.6\) for
B. The logit probabilities are approximately 44.1\% for A, 19.8\% for B,
and 36.1\% for the outside option.

Now introduce bundle C with the same modeled attributes as A. The logit
model assigns A and C approximately 30.6\% each, B 13.7\%, and the
outside option 25.1\%. C's share is 30.6 percentage points, but the
total purchasing share rises only from 63.9\% to 74.9\%, an increase of
about 11.0 points. The remaining share comes from the existing products.
Forecasting a new product's sales without accounting for substitution
would overstate incremental demand.

\begin{figure}[htbp]
\centering
% Original synthetic teaching figure; see corresponding module.
\begin{tikzpicture}[x=11cm,y=1cm]
\node[anchor=east] at (-.02,1.5) {Before launch};
\node[anchor=east] at (-.02,0) {After launch};
\fill[IlliniBlue] (0,1.1) rectangle (.4409,1.9);
\fill[Patina] (.4409,1.1) rectangle (.639,1.9);
\fill[LightStorm] (.639,1.1) rectangle (1,1.9);
\node[text=white] at (.22045,1.5) {A: 44.1\%};
\node[text=white] at (.54,1.5) {B: 19.8\%};
\node at (.8195,1.5) {Outside: 36.1\%};
\fill[IlliniBlue] (0,-.4) rectangle (.306, .4);
\fill[Patina] (.306,-.4) rectangle (.443483,.4);
\fill[IlliniOrange] (.443483,-.4) rectangle (.749475,.4);
\fill[LightStorm] (.749475,-.4) rectangle (1,.4);
\node[text=white] at (.153,0) {A: 30.6\%};
\node[text=white,font=\scriptsize,align=center] at (.3747,0) {B:\\13.7\%};
\node[text=IlliniBlue] at (.5965,0) {C: 30.6\%};
\node[align=center] at (.87475,0) {Outside:\\25.1\%};
\draw[IlliniBlue,dashed] (.639,.95)--(.639,-.85);
\draw[<->,very thick,IlliniOrange] (.639,-.75)--(.749475,-.75);
\node[align=center] at (.53,-1.55) {C gains 30.6 percentage points of share.\\Category share grows by only about 11.0 points.};
\end{tikzpicture}

\caption{A new product’s sales include substitution from existing products.}\label{fig:M04_cannibalization}
\end{figure}

With utility coefficients and the potential-customer population fixed,
adding C changes only the available choice set. Because C has the same
modeled attributes as A, it is a close substitute. The logit model's
independence-of-irrelevant-alternatives property makes C draw share
proportionally from all existing alternatives. Whether that pattern is
plausible is an empirical question. Richer substitution models may be
needed when close substitutes compete especially strongly with one
another.

A profit calculation must attach contribution margins and fixed costs to
the simulated quantities. If A and C each contribute \$4 per purchase
and B contributes \$3, expected contribution per potential customer
rises from about \$2.36 to \$2.86 before product-development and
operating costs. The contribution comparison is conditional on the
specified utilities, margins, and fixed customer population. A pilot or
holdout choice tasks can help evaluate the proposed product and model
before a larger decision.

\hypertarget{M04-estimate-and-validate-preferences}{%
\section{Estimate and Validate
Preferences}\label{M04-estimate-and-validate-preferences}}

\noindent A conjoint workflow begins with the managerial choice and feasible
attributes, then constructs the task design and recruits a relevant
respondent sample. Inspect completion time, straight-line responses,
inconsistent attention, and whether participants understood the
attributes. Exclusion rules should be set and reported transparently.
Removing respondents because their preferences are inconvenient changes
the target population.

Estimate a transparent baseline before adopting a more elaborate model.
Check coefficient coding, reference levels, price units, and whether
probabilities sum to one for every choice set. Evaluate prediction on
held-out tasks or respondents according to the intended use. A model
that fits familiar respondents' repeated choices may generalize
differently to a new customer population.

A useful simulation changes one decision at a time and records what is
fixed: attributes, price, competitor options, market size, and margins.
Share of preference from a survey is not automatically market share.
Distribution, awareness, availability, purchase frequency, and sampling
can create differences. State the link between the conjoint output and
the demand quantity used in the business calculation.

\hypertarget{M04-segmentation-begins-with-a-definition-of-similarity}{%
\section{Segmentation Begins With a Definition of
Similarity}\label{M04-segmentation-begins-with-a-definition-of-similarity}}

\noindent A retailer may want segments for differentiated service, product
positioning, or campaign design. Those objectives should guide the
variables used to define similarity. Grouping customers by recent
spending answers a different question from grouping them by category
preferences or service needs. A segment should be interpretable,
sufficiently stable, and connected to an action the firm can take.

Recency, frequency, and monetary value provide a useful descriptive
starting point. Define recency as time since the last eligible purchase
at a common observation date, frequency as the number of eligible
purchases during a fixed window, and monetary value as sales or
contribution during that window. Customers with different observation
histories require adjustment or a deliberate eligibility rule. A new
customer observed for one month should not be compared mechanically with
one observed for a year.

K-means selects \(K\) centers to minimize within-cluster squared
Euclidean distance:

\[
\sum_{i=1}^n\left\|z_i-\mu_{c(i)}\right\|^2.
\]

Here \(z_i\) is the vector of transformed customer features, \(c(i)\) is
customer \(i\)'s assigned cluster, and \(\mu_{c(i)}\) is its cluster
center. The feature transformations determine the distance.
Standardization often helps prevent a large numerical scale from
dominating, but it does not establish that all features should have
equal managerial importance.

For three customers with spending-frequency pairs A=(100,1), B=(110,9),
and C=(200,1), raw Euclidean distance makes B much closer to A than C
is. If spending is expressed in hundreds of dollars while frequency
remains in purchases, C becomes closer to A. The customer records have
not changed; the implicit weights have. This is why units and scaling
belong in the substantive definition of a segment.

\begin{figure}[htbp]
\centering
% Distances within each panel use the stated feature units.
\begin{tikzpicture}
\begin{axis}[name=raw,width=6.3cm,height=6.1cm,xmin=.5,xmax=2.5,ymin=0,ymax=120,ybar,bar width=18pt,xtick={1,2},xticklabels={A to B,A to C},ylabel={Euclidean distance},title={Raw spending and visits},nodes near coords,nodes near coords style={font=\scriptsize}]
\addplot[fill=IlliniBlue,draw=IlliniBlue] coordinates {(1,12.806)(2,100)};
\end{axis}
\begin{axis}[at={(raw.east)},anchor=west,xshift=2.3cm,width=6.3cm,height=6.1cm,xmin=.5,xmax=2.5,ymin=0,ymax=10,ybar,bar width=18pt,xtick={1,2},xticklabels={A to B,A to C},ylabel={Euclidean distance},title={Spending divided by 100},nodes near coords,nodes near coords style={font=\scriptsize}]
\addplot[fill=IlliniOrange,draw=IlliniOrange] coordinates {(1,8.0006)(2,1)};
\end{axis}
\end{tikzpicture}

\caption{Changing feature scales changes which customers appear similar.}\label{fig:M04_scaling}
\end{figure}

Changing spending units reverses the nearest neighbor because it changes
the relative weight placed on spending and visits. A segmentation report
should explain that metric and examine whether the proposed actions
remain sensible under reasonable transformations.

\hypertarget{M04-choosing-and-evaluating-segments}{%
\section{Choosing and Evaluating
Segments}\label{M04-choosing-and-evaluating-segments}}

\noindent Use more than one criterion to choose \(K\). Within-cluster dispersion
falls mechanically as more clusters are added. A silhouette measure
compares cohesion with separation under the chosen distance, but a high
silhouette does not establish business value. Examine segment sizes,
profiles, stability across resamples or time, and whether the segments
imply different feasible actions.

Hierarchical clustering provides a nested set of groupings whose
interpretation depends on the linkage and distance. Density-based
methods can identify irregular groups and mark some observations as
noise, with sensitivity to neighborhood parameters. Mixture models
provide probabilistic membership under distributional assumptions. These
alternatives address different patterns; no algorithm discovers a
uniquely correct segmentation independent of representation and purpose.

For customer activation, avoid using future purchase outcomes to define
a segment that supposedly exists before the campaign. For evaluating
segmentation, examine behavior not used mechanically to construct it. A
high-spending segment will have high spending by design if spending is a
clustering input. A more useful question is whether the grouping helps
predict future needs or improves the result of a tested service policy.

Clustering shopping trips on basket size, category composition, and
spending can help separate trip types before comparing channels. Fit the
representation without allowing the desired channel comparison to
determine the grouping. Examine profiles and sensitivity, because a
trip-classification algorithm can also create misleading comparisons if
its features or thresholds differ across channels.

\hypertarget{M04-customer-lifetime-value-depends-on-timing-and-margin}{%
\section{Customer Lifetime Value Depends on Timing and
Margin}\label{M04-customer-lifetime-value-depends-on-timing-and-margin}}

\noindent Customer lifetime value is the discounted expected contribution from a
relationship, under a stated time horizon and behavioral model. For a
newly acquired customer, a simple finite-horizon expression is

\[
\operatorname{CLV}=\sum_{t=1}^{T}\frac{\Pr(\text{active at }t)\,m_t}{(1+d)^t}-a.
\]

Here \(t\) indexes periods after acquisition, \(m_t\) is expected
contribution conditional on being active in period \(t\), \(d\) is the
discount rate per period, and \(a\) is acquisition cost paid at time
zero. Contribution should deduct relevant variable service and
fulfillment costs. The expression assumes the conditional margin and
activity probability have been defined consistently; more general models
can use the expected product of activity and customer-specific margin.

Assume a stationary subscription process: a currently active customer
survives each period with probability \(r\), earns a constant
contribution \(m\) at the end of a period only if still active, and has
an infinite horizon. With \(r/(1+d)<1\), the geometric series gives

\[
\operatorname{CLV}=\frac{mr}{1+d-r}-a.
\]

This convention excludes an immediate current-period contribution.
Including one would add that contribution and change the formula.
Monthly retention requires monthly margins and a monthly discount rate.
Combining monthly retention with an annual discount rate produces an
inconsistent calculation.

With monthly contribution \(m=60\), retention \(r=0.80\), monthly
discount rate \(d=0.01\), and acquisition cost \(a=50\), CLV is
approximately \$178.57. Increasing retention to 0.90 raises it to
approximately \$440.91. The gain is nonlinear because retention affects
many future payments. These values are conditional on constant retention
and margin; changes in retention or margin over time would change the
valuation.

\begin{figure}[htbp]
\centering
% Original synthetic teaching figure; see corresponding module.
\begin{tikzpicture}
\begin{axis}[xmin=.5,xmax=.95,ymin=0,ymax=1000,xlabel={Per-period retention probability},ylabel={Lifetime value after acquisition cost (dollars)},xtick={.5,.6,.7,.8,.9,.95},ytick={0,200,400,600,800,1000},ymajorgrids=true]
\addplot[IlliniBlue,domain=.5:.95,samples=150,very thick] {60*x/(1.01-x)-50};
\addplot[IlliniOrange,only marks,mark=*,mark size=3pt] coordinates {(.8,178.5714)(.9,440.9091)};
\node[anchor=north west] at (axis cs:.8,170) {$r=.80$: \$178.57};
\node[anchor=south east] at (axis cs:.89,480) {$r=.90$: \$440.91};
\node[align=left,anchor=north west] at (axis cs:.52,920) {$m=60,\ d=.01,\ a=50$\\Payments begin after first retention};
\end{axis}
\end{tikzpicture}

\caption{The value of retention is nonlinear under a stationary subscription model.}\label{fig:M04_clv}
\end{figure}

With monthly contribution, discounting, and acquisition cost fixed, a
one-point retention increase has greater value at a higher baseline
retention rate. Uncertainty about long-run retention can therefore
dominate a CLV forecast.

\hypertarget{M04-subscription-and-noncontractual-relationships}{%
\section{Subscription and Noncontractual
Relationships}\label{M04-subscription-and-noncontractual-relationships}}

\noindent In a subscription setting, cancellation is often observed. Survival
analysis can model time to cancellation and account for right censoring
when observation ends before cancellation. A hazard describes the
probability of exit during a period conditional on survival to that
period. Cohort effects, tenure dependence, and changes in service
quality can make constant retention implausible.

In a noncontractual retail setting, the last purchase does not reveal
whether a customer has permanently left. Models such as BG/NBD represent
purchase frequency and latent dropout under specific stochastic
assumptions. A separate spending model may estimate value per
transaction. Those assumptions require checking, including the treatment
of returns and any assumed relationship between purchase frequency and
monetary value. A customer-level count alone does not identify permanent
churn without a model.

Evaluate value forecasts on a holdout period with the same horizon used
in the decision. Compare aggregate and customer-level calibration, rank
ordering, and errors for new or infrequent customers. A model can rank
customers well while overpredicting total value. For an acquisition
budget, the aggregate monetary error may be more consequential than a
small improvement in ranking.

\hypertarget{M04-integrated-case-which-segment-should-receive-retention-support}{%
\section{Integrated Case: Which Segment Should Receive Retention
Support?}\label{M04-integrated-case-which-segment-should-receive-retention-support}}

\noindent Suppose segment H has high predicted CLV but little response to service
outreach, while segment M has moderate predicted CLV and a larger
response. Targeting only H can be inferior even if its value forecast is
accurate. The relevant quantity is the change in expected future
contribution caused by outreach, net of the outreach cost. CLV describes
a baseline relationship; a retention intervention changes that
relationship under a separate causal model.

A useful analysis would estimate or specify segment-specific
preferences, margins, and retention, then test the proposed offer.
Product design might attract a segment with expensive service needs; a
seemingly attractive acquisition price might therefore have low
contribution. Segment labels should help organize the decision without
replacing these calculations. A final recommendation should show which
assumption changes the preferred segment or product.

For practice, recompute the subscription value with \(r=0.85\) and
verify the result before acquisition cost is subtracted. Add an
immediate \$60 contribution and explain the change in timing. In the
conjoint example, remove the outside option and observe that simulated
purchase incidence becomes one by construction. In the clustering
example, compare similarity using category shares instead of spending
levels and explain which managerial question the new representation
serves.

\hypertarget{M04-ai-assisted-implementation-and-readings}{%
\section{AI-Assisted Implementation and
Readings}\label{M04-ai-assisted-implementation-and-readings}}

\noindent Ask a coding assistant to calculate choice probabilities from a supplied
profile table, check that they sum to one, and separate a new product's
share from incremental category demand. For clustering, require
documented transformations, segment profiles, and a sensitivity
comparison. For CLV, request an explicit payment timeline and verify the
closed-form calculation against a long finite sum. These checks catch
common mistakes that can survive apparently reasonable charts.

\href{https://sawtoothsoftware.com/academics/teaching-resources}{Sawtooth's
academic materials} provide conjoint examples.
\href{https://link.springer.com/book/10.1007/978-981-10-8818-6}{Dolnicar,
Grün, and Leisch's segmentation book} develops the connection between
segmentation design and use.
\href{https://www.pymc-marketing.io/en/latest/notebooks/clv/clv_quickstart.html}{PyMC-Marketing's
CLV documentation} offers implementation examples whose model
assumptions and software versions should be checked before adaptation.

\hypertarget{M05-m5.-text-social-media-and-customer-networks}{%
\chapter{Text, Social Media, and Customer
Networks}\label{M05-m5.-text-social-media-and-customer-networks}}

\noindent A retailer receives thousands of reviews and support messages each week.
The service manager needs to distinguish delivery delays, product
defects, and billing problems, then decide which issues deserve
attention. A count of negative words is an incomplete representation:
``the refund was fast, but the replacement failed'' contains different
evaluations of two aspects of the experience. The analytical task is to
turn language into a defensible measurement of a concept relevant to
action.

Text analysis also raises a sampling problem. Customers who post reviews
can differ from those who remain silent, and the platform can affect
what gets posted or displayed. Social-network analysis adds another
layer: observed connections and similar behavior can reflect influence,
common circumstances, or selection into relationships.

\hypertarget{M05-define-the-construct-and-the-document}{%
\section{Define the Construct and the
Document}\label{M05-define-the-construct-and-the-document}}

\noindent A construct is the concept the analysis intends to measure, such as
delivery dissatisfaction, perceived authenticity, or purchase intent. A
label is the operational category assigned to an observation. A document
might be a complete review, one sentence, a support ticket, or a
conversation. Those units should follow the construct. A review-level
sentiment label can obscure a positive service experience and a negative
product experience in the same review.

Write an annotation codebook before producing labels at scale. It should
define each category, inclusion and exclusion rules, treatment of
ambiguity, and examples near category boundaries. If a message can
contain several problems, a multi-label task may be appropriate. If the
decision requires routing each ticket to one team, define a primary
issue or a routing priority separately from the presence of every issue.

Metadata can be informative and potentially misleading. Ratings may help
validate a sentiment measure but can also leak the intended label into a
prediction exercise. Product category, language, channel, and date can
reveal domain differences. Remove or protect personal information before
sending text to an external service, and use approved data access. A
scraped review is not automatically free of personal or contractual
restrictions.

\hypertarget{M05-representing-language-with-counts-and-weights}{%
\section{Representing Language With Counts and
Weights}\label{M05-representing-language-with-counts-and-weights}}

\noindent A bag-of-words representation counts tokens while largely ignoring their
order. An \(n\)-gram representation retains local sequences, so ``not
good'' can differ from ``good.'' Tokenization, case normalization,
punctuation, and stop-word rules change the resulting features. Removing
negation words can destroy information needed for sentiment. Removing
product names may help focus on attributes or may erase the very object
whose complaints the manager needs to identify.

Term frequency-inverse document frequency gives greater weight to terms
that are frequent in a document but less common across the corpus. One
simple convention is

\[
\operatorname{tfidf}_{wd}=\operatorname{tf}_{wd}\log\left(\frac{N}{\operatorname{df}_w}\right),
\]

where \(\operatorname{tf}_{wd}\) is the frequency of word \(w\) in
document \(d\), \(N\) is the number of documents, and
\(\operatorname{df}_w\) is the number of documents containing the word.
Implementations vary in smoothing, normalization, and term-frequency
definitions. Record the convention rather than assuming a library's
output matches this exact expression.

If ``delivery'' appears in two of three documents, its
inverse-document-frequency term under this convention is
\(\log(3/2)\approx0.405\). A word appearing in every document receives
zero weight. That weighting can help discriminate documents, but a
common term may still be substantively important. Statistical
distinctiveness and managerial importance are different criteria.

Fit a learned vocabulary and IDF weights on training text when
evaluating prediction. Learning them from the test corpus exposes
information about the evaluation distribution. The leakage may appear
minor compared with using labels, yet it still changes the task. A
pipeline makes the training-only transformation explicit and
reproducible.

\hypertarget{M05-embeddings-and-semantic-similarity}{%
\section{Embeddings and Semantic
Similarity}\label{M05-embeddings-and-semantic-similarity}}

\noindent An embedding maps a document or token into a vector learned from
language data. Similar vectors can capture related meanings even when
documents use different words. Cosine similarity compares the angle
between vectors:

\[
\operatorname{cos}(u,v)=\frac{u^\top v}{\|u\|\,\|v\|}.
\]

The normalization removes vector length from the comparison. This is
useful when direction represents meaning in the selected model, but it
does not guarantee that the similarity corresponds to the marketing
construct. Two messages about refunds can be semantically close while
expressing opposite evaluations. Retrieval quality should be checked
with examples that matter for the downstream task.

A simple TF-IDF classifier can be a strong baseline for a well-defined
domain. Embeddings may help with paraphrases, sparse labels, and
semantic retrieval, but they introduce dependence on a pretrained model
and its training domain. Compare alternatives on the same held-out data
and operational criteria. Record model versions because an updated
representation can change nearest neighbors and labels even when the
input text is unchanged.

For interpretation, inspect representative successes and failures. A
model may associate a brand name with complaints because one product had
a recall during training. It can then misclassify neutral mentions after
the issue is resolved. Temporal validation and product-level holdouts
help reveal whether performance depends on such shortcuts.

\hypertarget{M05-sentiment-aspect-labels-and-topic-models}{%
\section{Sentiment, Aspect Labels, and Topic
Models}\label{M05-sentiment-aspect-labels-and-topic-models}}

\noindent A sentiment lexicon assigns scores to words and combines them using
rules. It offers a transparent baseline, but sarcasm, domain-specific
meanings, and context can cause errors. A supervised classifier learns a
mapping from labeled examples to outcomes. Its labels and training
sample define what it learns; a generic ``positive/negative'' dataset
may not represent customer-service severity.

Aspect-based analysis separates the object being evaluated from the
evaluation itself. For the sentence ``delivery was late, but the support
agent solved the problem,'' a useful record might contain
delivery-negative and service-positive labels. A single overall score
discards that distinction. If the action is a logistics intervention,
the aspect-specific measure is closer to the decision.

Topic models describe recurring patterns of word co-occurrence or
document similarity. A topic is an estimated pattern, not automatically
a stable consumer need. In probabilistic models, documents can combine
topics; in clustering approaches, assignments may be discrete or
probabilistic. Analysts supply topic names after inspecting high-weight
words and representative documents. Validate those names with actual
text and examine stability across preprocessing and model choices.

A manager should avoid interpreting an increase in a topic's share
without checking its denominator. More delivery complaints can mean more
delivery problems, more deliveries overall, a changed review prompt, or
fewer reviews about other topics. Combining text measures with
transaction or service-volume data helps distinguish these explanations.

\hypertarget{M05-case-study-a-complaint-detector-under-changing-prevalence}{%
\section{Case Study: A Complaint Detector Under Changing
Prevalence}\label{M05-case-study-a-complaint-detector-under-changing-prevalence}}

\noindent Suppose a classifier identifies a rare but costly product defect. In a
synthetic evaluation population of 1,000 tickets, 200 truly mention the
defect. The classifier detects 160 of them and incorrectly flags 80 of
the 800 other tickets. Sensitivity, or recall, is 80\%; specificity is
90\%; precision is 160/240, or 66.7\%.

\begin{table}[H]
\centering\small
\begin{tabular}{@{}lrrrr@{}}
\toprule
Population & True Defect Tickets & True Positives & False Positives &
Precision \\
\midrule

Evaluation sample & 200 & 160 & 80 & 66.7\% \\
Deployment sample & 10 & 8 & 99 & 7.5\% \\
\bottomrule
\end{tabular}
\end{table}

The second row holds sensitivity at 80\% and specificity at 90\%, but
reduces true prevalence to 1\%. Among 990 nondefect tickets, 99 are
flagged. Only eight of the 107 flagged tickets are true positives. A
manager using the original precision to plan review staffing would be
badly misled.

Bayes' rule makes the dependence explicit. If prevalence is \(\rho\),
sensitivity is \(s\), and specificity is \(q\), then

\[
\Pr(\text{defect}\mid\text{flag})=\frac{s\rho}{s\rho+(1-q)(1-\rho)}.
\]

The denominator counts both correct and incorrect flags. For a rare
issue, even a modest false-positive rate can dominate that total. This
is why validation should represent the deployment population or
explicitly adjust for a changed sampling design.

\begin{figure}[htbp]
\centering
% Original synthetic teaching figure; see corresponding module.
\begin{tikzpicture}
\begin{axis}[xmin=0,xmax=.3,ymin=0,ymax=1,xlabel={Prevalence of the target complaint},ylabel={Precision among flagged reviews},xtick={0,.1,.2,.3},ytick={0,.2,.4,.6,.8,1},ymajorgrids=true]
\addplot[IlliniBlue,very thick,domain=0:.3,samples=150] {.8*x/(.8*x+.1*(1-x))};
\addplot[IlliniOrange,only marks,mark=*,mark size=3pt] coordinates {(.01,.0747664)(.2,.6666667)};
\node[align=left,anchor=north west] at (axis cs:.12,.96) {Sensitivity fixed at 80\%\\Specificity fixed at 90\%};
\node[anchor=west] at (axis cs:.018,.13) {1\% prevalence: 7.5\% precision};
\node[anchor=north west] at (axis cs:.145,.58) {20\% prevalence: 66.7\%};
\end{axis}
\end{tikzpicture}

\caption{Precision falls as the target issue becomes rarer, even with unchanged sensitivity and specificity.}\label{fig:M05_prevalence}
\end{figure}

Sensitivity and specificity remain at 0.80 and 0.90 in both populations.
The decline in precision comes entirely from the lower prevalence.
Operational evaluation should also consider the cost of missed defects,
the burden of false alarms, and whether uncertain cases can be routed to
human review.

\hypertarget{M05-human-and-llm-annotation}{%
\section{Human and LLM Annotation}\label{M05-human-and-llm-annotation}}

\noindent LLMs can help draft a codebook, suggest candidate labels, or classify
text using explicit instructions. For a defensible measurement exercise,
preserve the original text, the codebook version, model identifier,
prompt, output schema, and date. Require a bounded output format with an
allowed ``uncertain'' category where appropriate. Validate the output
against the schema and inspect cases that violate the instructions.

A prompt should identify the construct, unit, categories, boundary
rules, and examples. For delivery dissatisfaction, it should distinguish
a complaint about an actual delay from a question about the promised
delivery date. Ask for a short supporting span from the supplied text
when that improves auditability. A fluent explanation is not evidence
that the label is correct; compare it with independent human judgments
on a relevant sample.

Human agreement is informative because the construct may itself be
ambiguous. Double-code a sample, examine disagreements, and refine
definitions before scaling. Report category-specific errors, prevalence,
and performance across languages or product groups. An average accuracy
can conceal systematic failure on the minority category that motivates
the analysis.

Treat text being classified as data. A review may contain language
resembling an instruction to the model. It should not change the
annotation task, request external actions, or supply new system rules.
Separate the fixed task instructions from the document content,
constrain outputs, and inspect unexpected behavior. Use synthetic or
approved text for classroom exercises involving external services.

\hypertarget{M05-networks-connections-reach-and-brokerage}{%
\section{Networks: Connections, Reach, and
Brokerage}\label{M05-networks-connections-reach-and-brokerage}}

\noindent A network represents entities as nodes and relationships as edges. Nodes
might be customers, creators, brands, or products. An edge might
indicate following, communication, co-purchase, or a similarity
criterion. Directed and weighted networks preserve different information
from undirected binary networks. The edge definition should follow the
substantive relationship the analyst intends to study.

Degree counts a node's direct connections. Betweenness summarizes how
often a node lies on shortest paths between other nodes. A person with
few direct contacts can have high betweenness if those contacts connect
otherwise separate groups. Eigenvector-based measures give more weight
to connections with well-connected nodes. These metrics describe network
structure under their definitions; they do not by themselves establish
persuasion, credibility, or the effect of paying a creator.

\begin{figure}[htbp]
\centering
% Original synthetic teaching figure; see corresponding module.
\begin{tikzpicture}[v/.style={circle,draw=IlliniBlue,fill=white,minimum size=9mm},b/.style={circle,draw=IlliniOrange,very thick,fill=white,minimum size=10mm}]
\node[v] (a) at (0,1.2) {A};\node[v] (b) at (0,-1.2) {B};\node[v] (c) at (1.5,0) {C};\node[v] (d) at (-1.5,0) {D};
\node[b] (r) at (4,0) {R};
\node[v] (e) at (6.5,0) {E};\node[v] (f) at (8,1.2) {F};\node[v] (g) at (8,-1.2) {G};\node[v] (h) at (9.5,0) {H};
\foreach \x/\y in {a/b,a/c,a/d,b/c,b/d,c/d,e/f,e/g,e/h,f/g,f/h,g/h} {\draw[IlliniBlue] (\x)--(\y);}
\draw[IlliniOrange,very thick] (c)--(r)--(e);
\node[align=center] at (0,-2.2) {Four connected customers};
\node[align=center] at (8,-2.2) {Four connected customers};
\node[align=center,text width=4cm] at (4,2.3) {R has only two links,\\but connects the groups.};
\draw[accent] (4,1.6)--(r.north);
\end{tikzpicture}

\caption{A broker with few direct links can connect two otherwise separate customer groups.}\label{fig:M05_network_bridge}
\end{figure}

Broker R has two direct connections, while each within-group connector
has four. Removing R disconnects the two groups. A customer with fewer
direct contacts can therefore be essential for reaching a different
audience. A marketing campaign may value either role, depending on
message transmission, audience overlap, and the intended outcome.

Observed similarity among connected customers can arise from influence,
homophily, or common exposure. Friends may adopt a product because they
affect one another, because they already share tastes, or because they
see the same promotion. Regressing adoption on friends' adoption does
not distinguish those mechanisms without additional design. Randomized
encouragement, timing restrictions, or credible natural variation can
help, with attention to spillovers and the unit of assignment.

\hypertarget{M05-social-media-case-engagement-and-commercial-value}{%
\section{Social-Media Case: Engagement and Commercial
Value}\label{M05-social-media-case-engagement-and-commercial-value}}

\noindent A brand compares two creators. One receives many comments, but those
comments largely come from an audience that already follows the brand.
The other reaches a smaller group with less overlap. A high engagement
count does not establish more incremental sales. The comparison requires
a target population, exposure definition, outcome horizon, and a way to
measure what would have happened without the collaboration.

Text can help describe comment content, and network overlap can describe
audience duplication. Those measurements support campaign design, but
they should be combined with an evaluation strategy. A test might
randomize eligible regions or audience groups when feasible, or compare
outcomes using a credible observational design. If communication crosses
groups, spillovers can complicate the interpretation of an untreated
group.

For classroom practice, construct two small audience sets and calculate
their overlap. Then compare comments per exposed user with total
comments. Explain why neither is a direct estimate of incremental
profit.

\hypertarget{M05-ai-assisted-workflow-practice-and-readings}{%
\section{AI-Assisted Workflow, Practice, and
Readings}\label{M05-ai-assisted-workflow-practice-and-readings}}

\noindent Start a text exercise with a small manually inspected sample and a
codebook. Build a transparent baseline, reserve an evaluation sample,
and compare a more complex method on the same outcomes. Ask an assistant
to produce a confusion matrix and representative errors, then
independently recompute precision and recall from the counts. Keep model
selection and prompt refinement away from the final test sample.

For practice, calculate accuracy in the first complaint-detector sample:
88\%. Explain why a classifier that always predicts ``no defect'' would
achieve 99\% accuracy in the rare-prevalence deployment sample while
detecting no defects. Next identify how removing ``not'' changes the
meaning of ``not reliable.'' Finally, explain why two highly connected
adopters may share preferences even if neither caused the other's
adoption.

The
\href{https://scikit-learn.org/stable/modules/feature_extraction.html\#text-feature-extraction}{scikit-learn
text-feature documentation} supports count and TF-IDF representations.
\href{https://github.com/cjhutto/vaderSentiment}{VADER's author
repository} provides a rule-based sentiment baseline.
\href{https://doi.org/10.1007/s11002-020-09535-7}{Hagen and colleagues
(2020)} discuss ways machine learning can support behavioral marketing
research.

\hypertarget{M06-m6.-causal-effects-and-marketing-experiments}{%
\chapter{Causal Effects and Marketing
Experiments}\label{M06-m6.-causal-effects-and-marketing-experiments}}

\noindent A retailer wants to know whether an email campaign increases
contribution. Customers who open marketing emails often purchase more,
but they may already be more engaged. Comparing openers with nonopeners
can therefore reward a campaign for behavior it did not cause. The
decision requires a comparison between outcomes under alternative
actions for the same target population.

\hypertarget{M06-define-the-treatment-outcome-and-estimand}{%
\section{Define the Treatment, Outcome, and
Estimand}\label{M06-define-the-treatment-outcome-and-estimand}}

\noindent Let \(Y_i(1)\) be customer \(i\)'s 30-day contribution if assigned an
email and \(Y_i(0)\) the contribution if not assigned. The individual
effect is \(Y_i(1)-Y_i(0)\). Only one of these outcomes is observed for
a given customer under a single assignment. A causal design supplies a
defensible way to learn an average comparison despite that missing
counterfactual.

For a target population, the average treatment effect is

\[
\operatorname{ATE}=E[Y_i(1)-Y_i(0)].
\]

The expectation refers to the stated population and outcome window. An
average treatment effect on the treated uses the population that
actually received treatment. An intention-to-treat effect uses
assignment as the treatment, even if some assigned customers never open
the email. These are different estimands. The managerial decision to
send an email often corresponds closely to the assignment effect,
because opening is outside the firm's direct control.

A well-defined treatment includes the offer, timing, channel, and
implementation. ``Digital marketing'' is too broad to identify a single
comparison. The outcome should include economically relevant costs and
have enough follow-up to capture returns or delayed purchases. If
customers can affect one another, an individual's outcome may depend on
other customers' assignments. The unit and design then need to account
for interference rather than silently assuming it away.

\hypertarget{M06-why-the-treated-untreated-difference-can-be-biased}{%
\section{Why the Treated-Untreated Difference Can Be
Biased}\label{M06-why-the-treated-untreated-difference-can-be-biased}}

\noindent The observed difference between treated and untreated customers can be
decomposed as

\[
E[Y\mid D=1]-E[Y\mid D=0]
=E[Y(1)-Y(0)\mid D=1]
+\{E[Y(0)\mid D=1]-E[Y(0)\mid D=0]\}.
\]

Here \(D\) indicates treatment. The first term is the average treatment
effect among treated customers. The bracketed term compares their
untreated outcomes with untreated customers' outcomes. If the firm
targets people who would purchase anyway, that selection term can be
positive. If it targets customers in serious decline, the term can be
negative. The direction follows the assignment process, not a universal
rule about observational data.

An unobserved engagement variable may affect both email opening and
purchasing. Adjusting for pre-treatment engagement measures can be
useful if they capture the relevant confounding, but controlling for a
variable does not guarantee this condition. A post-treatment variable
such as clicking the email can mediate part of its effect. Conditioning
on customers who click changes the population and can introduce
selection.

\begin{figure}[htbp]
\centering
% Original synthetic teaching figure; see corresponding module.
\begin{tikzpicture}
\node[box,text width=2.8cm] (z) at (0,0) {Randomized\\assignment $Z$};
\node[orangebox,text width=2.8cm] (o) at (5,0) {Email opening $O$};
\node[box,text width=2.8cm] (y) at (10,0) {Purchase $Y$};
\node[box,text width=3cm] (u) at (7.5,2.5) {Prior\\engagement $U$};
\draw[edge] (z)--(o);\draw[edge] (o)--(y);\draw[edge] (u)--(o);\draw[edge] (u)--(y);
\draw[edge] (z.south)--(0,-1.4)--node[below,font=\scriptsize] {Other pathways from assignment} (10,-1.4)--(y.south);
\node[align=center,text width=11cm] at (5,-2.9) {Opening depends on engagement and assignment.\\Keep all assigned customers in the analysis.};
\end{tikzpicture}

\caption{Assignment, engagement, and opening define different causal comparisons.}\label{fig:M06_dag}
\end{figure}

The firm randomizes email assignment, while customers choose whether to
open. Assignment can affect opening, and engagement can affect both
opening and purchases. Restricting attention to openers creates a
selected group. An assignment-based analysis preserves the randomized
comparison when its implementation and outcome measurement are valid.

\hypertarget{M06-designing-an-experiment-around-the-decision}{%
\section{Designing an Experiment Around the
Decision}\label{M06-designing-an-experiment-around-the-decision}}

\noindent An experiment needs an eligible population, assignment unit, treatment
implementation, primary outcome, horizon, and analysis plan. Randomize
at a unit that corresponds to exposure and likely spillovers. A
customer-level email experiment may be appropriate when accounts act
independently. A store-level display intervention requires store
assignment or a design that accounts for shared exposure. Household,
geographic, or network spillovers can require broader clusters.

Balance treatment groups through randomization, with blocking or
stratification when important pre-treatment characteristics are known.
For example, randomize within customer-tenure groups to improve balance
and precision. Record assignment before contact and retain every
assigned unit in the primary analysis when estimating the
intention-to-treat effect. Excluding nonopeners, customers who later
unsubscribe, or stores with disappointing execution can undermine the
randomized comparison.

Specify a primary outcome and guardrails before viewing results.
Incremental contribution might be primary, with unsubscribe rates,
complaints, and returns as guardrails. Multiple outcomes can be
informative, but selecting the most favorable result after examining
many tests exaggerates evidence. Decide how multiple testing, sequential
monitoring, and stopping will be handled. Ordinary fixed-horizon
p-values do not generally retain their advertised interpretation under
unrestricted repeated peeking and outcome selection.

\hypertarget{M06-case-study-estimate-an-email-assignment-effect}{%
\section{Case Study: Estimate an Email Assignment
Effect}\label{M06-case-study-estimate-an-email-assignment-effect}}

\noindent In a synthetic experiment, 1,000 customers are assigned an email and
1,000 are assigned no email. Within 30 days, 130 assigned customers and
100 controls make at least one eligible purchase. The estimated
assignment effect on purchase incidence is

\[
\hat\tau=\frac{130}{1000}-\frac{100}{1000}=0.03,
\]

or three percentage points. Relative to the control rate of 10\%, this
is a 30\% relative increase. The two descriptions use different
denominators; reporting only ``30\% lift'' can obscure the absolute
change.

Under independent customer assignments and the usual large-sample
approximation, an unpooled standard error is

\[
\widehat{\operatorname{SE}}(\hat\tau)
=\sqrt{\frac{0.13(0.87)}{1000}+\frac{0.10(0.90)}{1000}}
\approx0.0143.
\]

An approximate 95\% interval is {[}0.0021, 0.0579{]}, or about 0.21 to
5.79 percentage points. This interval describes sampling uncertainty
under the stated design and approximation. It does not account for an
unmeasured outcome-definition error or interference between assigned
customers.

Suppose an additional purchase contributes \$40 and email delivery costs
\$0.50 per assigned customer. If the offer has no additional discount
cost and purchase incidence is the relevant incremental quantity,
estimated incremental contribution is \$40(0.03)-\$0.50=\$0.70 per
assigned customer. Applying the same transformation to the effect
interval gives approximately \$-0.42 to \$1.82. A positive point
estimate therefore coexists with economically consequential uncertainty.
A decision can combine that uncertainty with rollout costs, the value of
another test, and guardrail outcomes.

The margin assumption matters. If treatment also changes basket size,
returns, or discount expense, purchase incidence alone is insufficient.
An experiment measuring contribution directly may better match the
objective.

\hypertarget{M06-power-precision-and-the-minimum-useful-effect}{%
\section{Power, Precision, and the Minimum Useful
Effect}\label{M06-power-precision-and-the-minimum-useful-effect}}

\noindent Power planning begins with the smallest effect that could change the
decision, not with a convenient number of rows. For two equally sized
independent groups and a binary outcome with baseline rate near \(p\), a
rough per-group sample-size approximation is

\[
n\approx\frac{2p(1-p)(z_{1-\alpha/2}+z_{1-\beta})^2}{\delta^2},
\]

where \(\delta\) is the absolute effect in probability units, \(\alpha\)
is the two-sided significance level, and \(1-\beta\) is target power.
This approximation uses a common baseline variance and normal critical
values. More precise planning can use the expected rates under the
alternative and the actual assignment design.

With \(p=0.10\), \(\delta=0.02\), \(\alpha=0.05\), and 80\% power, using
1.96 and 0.84 gives about 3,528 customers per group. Halving the target
effect approximately quadruples the required sample under this formula.
A two-percentage-point effect is 0.02, not 2; a unit error can produce a
meaningless sample-size recommendation.

Clustering reduces the amount of independent information. A common
approximation with equal cluster size \(m\) and intracluster correlation
\(\rho\) multiplies individual-level variance by \(1+(m-1)\rho\). For
\(m=20\) and \(\rho=0.05\), that factor is 1.95. The number of clusters,
unequal cluster sizes, assignment design, and small-sample inference
still require explicit consideration. Thousands of customers across four
stores do not provide thousands of independent store-level treatment
assignments.

\hypertarget{M06-implementation-checks-that-protect-the-comparison}{%
\section{Implementation Checks That Protect the
Comparison}\label{M06-implementation-checks-that-protect-the-comparison}}

\noindent Verify that the assignment counts match the intended allocation and
investigate material sample-ratio discrepancies. Check whether tracking
differs between treatment and control, whether treatment changes the
probability an outcome is observed, and whether exclusions were applied
consistently. Pre-treatment balance is a diagnostic, not proof that
every causal assumption holds. In a correctly randomized experiment,
chance imbalances can occur.

Noncompliance changes the estimand. If only 40\% of assigned customers
open the email, the intention-to-treat effect remains the effect of
assignment under the observed implementation. Dividing it by an
opening-rate difference identifies a complier effect only under
additional instrumental-variable conditions, including an appropriate
exclusion restriction and monotonicity. The firm may care about the
assignment effect even when a different scientific question concerns
actual exposure.

Use pre-treatment outcomes to improve precision when justified. A
regression adjustment can compare outcomes while accounting for baseline
spending or stratification variables. Avoid selecting an adjustment
because it makes the treatment coefficient significant. Report the
planned analysis and explain any departures. Analyze heterogeneity
cautiously, especially when many subgroups are examined after seeing the
main result.

\hypertarget{M06-observational-adjustment-and-its-limits}{%
\section{Observational Adjustment and Its
Limits}\label{M06-observational-adjustment-and-its-limits}}

\noindent When randomization is unavailable, a common approach assumes that
treatment is conditionally independent of potential outcomes given
observed pre-treatment covariates, with overlap in treatment
probabilities. Regression adjustment, matching, and inverse-probability
weighting implement different ways to construct a comparison under that
condition. The central question is whether the observed covariates make
the treated and untreated outcomes comparable.

A propensity score is the treatment probability given those covariates.
Extreme estimated scores indicate limited overlap and can create
unstable weights. Trimming or restricting the sample changes the
population for which the effect is estimated. Report that population
explicitly. A flexible prediction algorithm for treatment assignment
does not establish that all confounders were measured.

A doubly robust estimator combines an outcome model and a treatment
model and can be consistent when one nuisance model is correctly
specified, under the underlying identification conditions and regularity
assumptions. This robustness does not cover unmeasured confounding. For
advanced implementations, sample splitting and cross-fitting help
control overfitting of nuisance functions; the design remains the
foundation of the causal interpretation.

\hypertarget{M06-difference-in-differences-and-a-counterfactual-trend}{%
\section{Difference-in-Differences and a Counterfactual
Trend}\label{M06-difference-in-differences-and-a-counterfactual-trend}}

\noindent A retailer introduces an app in treated markets while comparable markets
remain untreated. A before-after difference in treated markets mixes the
app's effect with common demand changes. Difference-in-differences
subtracts the comparison group's change:

\[
\hat\tau_{\mathrm{DiD}}=(\bar Y_{T,post}-\bar Y_{T,pre})-(\bar Y_{C,post}-\bar Y_{C,pre}).
\]

In a synthetic example, treated-market sales rise from 120 to 150 and
comparison-market sales rise from 100 to 110. The DiD estimate is 20
sales units. Under the identifying parallel-trends condition, the
treated markets' untreated counterfactual would have risen to 130.
Different starting levels are compatible with this design; the condition
concerns untreated changes.

\begin{figure}[htbp]
\centering
% Original synthetic teaching figure; see corresponding module.
\begin{tikzpicture}
\begin{axis}[xmin=-.12,xmax=1.38,ymin=90,ymax=162,xtick={0,1},xticklabels={Before,After},ylabel={Outcome (illustrative units)},legend style={at={(.5,1.08)},anchor=south,legend columns=3}]
\addplot[IlliniOrange,mark=*,very thick] coordinates {(0,120)(1,150)};
\addplot[IlliniBlue,mark=square*,very thick] coordinates {(0,100)(1,110)};
\addplot[Storm,dashed,mark=o,very thick] coordinates {(0,120)(1,130)};
\legend{Treated,Comparison,Treated counterfactual}
\draw[<->,IlliniOrange,very thick] (axis cs:1.08,130)--(axis cs:1.08,150);
\node[anchor=west,align=left] at (axis cs:1.1,140) {Effect\\$=20$};
\end{axis}
\end{tikzpicture}

\caption{Difference-in-differences subtracts the comparison group’s change to construct the counterfactual.}\label{fig:M06_did}
\end{figure}

The dashed line is a counterfactual implied by parallel untreated
trends, not an observed sales series. The vertical gap between observed
treated sales and that line is the estimated effect. A
treatment-specific local shock at the same time would challenge the
interpretation even if the pre-treatment levels looked similar.

Inspect multiple pre-treatment periods, institutional timing, possible
anticipation, spillovers, and changes in sample composition. Similar
pre-trends can support the design's plausibility, but a nonsignificant
pre-trend test does not establish parallel counterfactual trends. With
staggered adoption and heterogeneous effects, conventional two-way
fixed-effects coefficients can combine comparisons in problematic ways.
Use estimators aligned with cohort and time-specific effects, with a
clear comparison group, when extending beyond the simple two-period
design.

\hypertarget{M06-other-designs-discontinuities-and-instruments}{%
\section{Other Designs: Discontinuities and
Instruments}\label{M06-other-designs-discontinuities-and-instruments}}

\noindent A regression discontinuity design can use a treatment rule at a
threshold, such as eligibility based on a pre-existing score. Customers
just above and below the cutoff may be comparable if other determinants
of the outcome vary smoothly there and the score is not precisely
manipulated. The resulting effect is local to the cutoff and the
relevant compliance behavior. Bandwidth, functional form, density
patterns, and discontinuities in other predetermined variables deserve
inspection.

An instrumental variable changes treatment while affecting the outcome
only through the treatment channel, under the design's conditions.
Relevance is an empirical first-stage requirement; exclusion is a
substantive condition. A strong first stage cannot establish exclusion.
In a binary-treatment setting with independence, exclusion, and
monotonicity, the familiar local average treatment effect concerns
compliers whose treatment changes with the instrument. It need not equal
the population ATE.

Marketing instruments require particularly careful reasoning because
pricing, distribution, and exposure decisions often respond to demand. A
variable correlated with treatment is not sufficient. Explain who makes
the relevant decision, when the instrument is determined, and which
other channels could connect it to outcomes. If those channels remain
credible, the causal interpretation needs revision.

\hypertarget{M06-spillovers-and-experiment-units}{%
\section{Spillovers and Experiment
Units}\label{M06-spillovers-and-experiment-units}}

\noindent An advertising campaign can shift purchases between stores or products
rather than increase total demand. A customer-level outcome may show a
gain while chain-wide contribution changes little. Define the outcome at
the level relevant to the managerial objective and examine displacement
where it is plausible. The same issue arises when customers share offers
or when a platform changes recommendations for multiple users.

\begin{figure}[htbp]
\centering
% Original synthetic teaching figure; see corresponding module.
\begin{tikzpicture}[v/.style={circle,draw=IlliniBlue,minimum size=11mm},t/.style={v,fill=IlliniOrange}]
\node[t] (t) at (0,0) {T};\node[v] (a) at (2,1.5) {U};\node[v] (b) at (2,-1.5) {U};\node[v] (c) at (4,0) {U};
\draw[accent] (t)--(a);\draw[accent] (t)--(b);\draw[edge,dashed] (a)--(c);\draw[edge,dashed] (b)--(c);
\draw[IlliniBlue,dashed,rounded corners] (-1,-2.2) rectangle (5,2.2);
\node[align=center] at (2,-3) {One connected household or market\\can form a randomization cluster.};
\node[box,text width=4cm] at (8,1) {T: assigned treatment\\U: assigned control};
\node[align=left,text width=4cm] at (8,-1.1) {Control status does not guarantee absence of exposure through other people.};
\end{tikzpicture}

\caption{A treatment can affect both assigned units and nearby untreated units.}\label{fig:M06_spillovers}
\end{figure}

An assigned store can gain sales while diverting purchases from another
store. If the comparison store loses customers, subtracting its decline
can overstate the increase in total chain sales. Cluster randomization,
geographic separation, or designs that vary treatment saturation may
help identify the relevant effects, subject to their own practical
constraints.

\hypertarget{M06-ai-assisted-workflow-and-practice}{%
\section{AI-Assisted Workflow and
Practice}\label{M06-ai-assisted-workflow-and-practice}}

\noindent Ask an assistant to state the estimand, assignment unit, outcome
horizon, and comparison group before writing estimation code. For the
email example, require the conversion effect in percentage points, the
unpooled standard error, and the contribution calculation. Check the
arithmetic independently. Ask it to explain why restricting the sample
to email openers changes the analysis.

For observational data, draw or describe the assignment process before
selecting controls. Request an overlap diagnostic and a report of which
observations receive large weights. For DiD, ask the assistant to show
observed and counterfactual series separately and identify the
assumptions that connect them. A convincing plot should not substitute
for those conditions.

If comparison-market sales rise by 25 rather than 10, the DiD estimate
becomes five. If email cost doubles to \$1, the point estimate of
incremental contribution falls to \$0.20. If treatment raises returns,
explain why a conversion effect can overstate the contribution effect.
If two stores share most customers, identify why store assignment may
create interference.

\hypertarget{M06-readings}{%
\section{Readings}\label{M06-readings}}

\noindent \href{https://matheusfacure.github.io/python-causality-handbook/landing-page.html}{Facure's
causal-inference handbook} supplies accessible implementations of
several designs.
\href{https://exp-platform.com/Documents/GuideControlledExperiments.pdf}{Kohavi
and colleagues' practical guide} discusses controlled online
experiments.
\href{https://doi.org/10.1016/j.jeconom.2020.12.001}{Callaway and
Sant'Anna (2021)} develop difference-in-differences with multiple
periods.

\hypertarget{M07-m7.-prediction-machine-learning-and-marketing-actions}{%
\chapter{Prediction, Machine Learning, and Marketing
Actions}\label{M07-m7.-prediction-machine-learning-and-marketing-actions}}

\noindent A retailer wants to anticipate next month's customer purchases and
allocate limited service capacity. A model that predicts historical
purchases accurately may still fail when deployed: its features may
contain future information, its test set may reuse the same customers as
training, or its objective may reward errors that are costly for the
business. Prediction should therefore be designed around an explicit
deployment situation.

\hypertarget{M07-specify-the-forecast-origin-and-target}{%
\section{Specify the Forecast Origin and
Target}\label{M07-specify-the-forecast-origin-and-target}}

\noindent A prediction task needs a target variable, a forecast origin, a horizon,
and an eligible population. For customer churn, define what counts as
churn and when that status becomes observable. In a subscription
business, cancellation may be explicit. In noncontractual retail, ``no
purchase for 90 days'' is an operational label, not proof that a
relationship has ended permanently.

Construct predictors using information available at the forecast origin.
If predicting purchases during November, features might use transactions
through October 31. November support contacts, returns, or campaign
interactions are unavailable at that origin. A database may contain them
when the analyst extracts the file, making leakage easy to overlook.
Store feature timestamps or document their measurement windows.

The deployment population determines the validation design. Predicting
future behavior for existing customers differs from predicting behavior
for newly acquired customers. A random row split can place the same
customer in training and test data, giving a favorable assessment for a
task that actually requires generalization to new customers. A time
split addresses future periods, while a customer-group split addresses
new customers. Some tasks require both.

\begin{figure}[htbp]
\centering
% Original synthetic teaching figure; see corresponding module.
\begin{tikzpicture}[x=1.55cm,y=.85cm]
\foreach \x/\lab in {0/Period 1,1/Period 2,2/Period 3} {\node at (\x,1) {\lab};\node at (\x+4.5,1) {\lab};}
\node[font=\small\bfseries] at (1,2) {Random rows};\node[font=\small\bfseries] at (5.5,2) {Future-period holdout};
\foreach \y/\lab in {0/A,-1/B,-2/C} {\node[anchor=east] at (-.55,\y) {\lab};\node[anchor=east] at (3.95,\y) {\lab};}
\foreach \x/\y in {0/0,2/0,1/-1,2/-1,0/-2,1/-2} {\fill[IlliniBlue] (\x-.43,\y-.36) rectangle (\x+.43,\y+.36);}
\foreach \x/\y in {1/0,0/-1,2/-2} {\fill[IlliniOrange] (\x-.43,\y-.36) rectangle (\x+.43,\y+.36);}
\foreach \x in {4.5,5.5} {\foreach \y in {0,-1,-2} {\fill[IlliniBlue] (\x-.43,\y-.36) rectangle (\x+.43,\y+.36);}}
\foreach \y in {0,-1,-2} {\fill[IlliniOrange] (6.07,\y-.36) rectangle (6.93,\y+.36);}
\node[align=center,text width=12cm] at (3.1,-3.4) {Blue: training. Orange: test. Rows identify customers.\\For new-customer deployment, hold out complete customer rows instead.};
\end{tikzpicture}

\caption{The validation split should reproduce the information available at deployment.}\label{fig:M07_splits}
\end{figure}

A random row split can reuse a customer's history in training and test
data. Holding out complete customers evaluates generalization to new
customers; holding out later periods evaluates future performance. The
intended use determines which comparison is relevant.

\hypertarget{M07-baselines-loss-functions-and-regularization}{%
\section{Baselines, Loss Functions, and
Regularization}\label{M07-baselines-loss-functions-and-regularization}}

\noindent A baseline should be operationally meaningful. For sales forecasting,
compare with the previous period, the same period last year, or a simple
customer-history average when those quantities are available. For
classification, compare with a constant probability and a transparent
logistic regression. A complex model's incremental value is easier to
assess when these comparisons are reported on the same test
observations.

A supervised learner minimizes a loss that quantifies prediction error.
Squared error penalizes large deviations strongly; absolute error gives
a different compromise under skewness. For a binary outcome, log loss is

\[
L=-\frac{1}{n}\sum_{i=1}^n\{y_i\log\hat p_i+(1-y_i)\log(1-\hat p_i)\}.
\]

Here \(y_i\) is the observed binary outcome and \(\hat p_i\) is its
predicted probability. Confident mistakes receive large penalties. The
loss evaluates probability forecasts; it does not by itself encode the
cost of a marketing action. Select a loss appropriate to the prediction
task, then evaluate how the forecasts support the downstream decision.

Regularization constrains a flexible model to reduce overfitting. Ridge
regression adds a squared-coefficient penalty, while lasso adds an
absolute-coefficient penalty and can set some coefficients to zero.
Standardize features appropriately when the penalty depends on
coefficient scale, and fit that transformation using training data. The
regularization strength is a tuning choice selected with validation
data, not with the final test sample.

Bias and variance describe an important tradeoff. A rigid model may
systematically miss a pattern; a very flexible model may fit noise
specific to the training sample. More data can help, but only if it
represents the deployment problem and adds relevant information. A
million duplicated or highly dependent rows can provide less useful
information than their count suggests.

\hypertarget{M07-logistic-regression-trees-and-ensembles}{%
\section{Logistic Regression, Trees, and
Ensembles}\label{M07-logistic-regression-trees-and-ensembles}}

\noindent Logistic regression maps a linear index to a probability:

\[
\hat p_i=\frac{1}{1+\exp[-(\beta_0+x_i^\top\beta)]}.
\]

The coefficients describe changes in log odds under the model. A
coefficient's effect on probability depends on the starting probability
and other predictors. Predicted probabilities are often easier to
communicate than odds ratios, especially when the business decision uses
expected counts or contribution. Nonlinear features and interactions can
extend the model while preserving a relatively transparent structure.

A decision tree partitions the feature space into regions using
sequential splits. It can represent interactions without specifying them
manually: a spending threshold might matter only among recent customers.
Deep trees can fit small idiosyncratic groups, making pruning, minimum
leaf sizes, and validation important. A split point near an observed
value should not automatically be interpreted as a behavioral threshold.

Random forests average trees fitted with randomness in observations and
candidate features, reducing instability relative to a single tree in
many settings. Gradient boosting builds an ensemble sequentially to
improve a chosen loss, with tuning choices such as tree depth, learning
rate, and number of iterations. These methods can be effective on
tabular data, but performance remains empirical. More complexity also
increases the burden of checking leakage, calibration, computation, and
reproducibility.

Neural networks learn layered representations and can be especially
useful for text, images, or large complex datasets. Their flexibility
does not eliminate the need for baselines or careful evaluation. In a
modest customer table, a well-tuned simpler model may be competitive and
easier to maintain. Choose the method based on task performance and
practical constraints rather than the apparent modernity of the
algorithm.

\hypertarget{M07-validation-cross-validation-and-leakage}{%
\section{Validation, Cross-Validation, and
Leakage}\label{M07-validation-cross-validation-and-leakage}}

\noindent Use training data to fit parameters, validation data or cross-validation
to select models and tuning choices, and a final test sample to estimate
performance after those choices are fixed. Reusing the test sample
repeatedly turns it into a validation sample, even if the file is still
named ``test.'' Report the procedure so the claimed holdout remains
meaningful.

All learned preprocessing belongs inside the training process. This
includes imputation, scaling, feature selection, vocabulary
construction, dimensionality reduction, and target-based encodings. In
cross-validation, fit each transformation on the training fold and apply
it to that fold's validation observations. A pipeline makes the sequence
explicit. The
\href{https://scikit-learn.org/stable/common_pitfalls.html}{scikit-learn
pitfalls guide} documents implementation patterns for this separation.

Time-series and grouped data need adapted splits. Rolling-origin
evaluation trains on earlier periods and predicts later ones,
potentially repeating this across several origins. Grouped validation
holds out complete customers, households, stores, or other relevant
units. If a target window overlaps a training outcome window, a temporal
gap may be necessary. Document the dates and units assigned to each
split.

A useful audit asks when each feature became observable. In a
marketing-call dataset, call duration can be useful after a call but
unavailable when deciding whom to call. A model using that feature
answers a different prediction problem. The feature's legitimacy follows
the decision time, not whether a package accepts it or whether it
improves accuracy.

\hypertarget{M07-case-study-accuracy-can-conceal-a-useless-model}{%
\section{Case Study: Accuracy Can Conceal a Useless
Model}\label{M07-case-study-accuracy-can-conceal-a-useless-model}}

\noindent Suppose 5\% of 1,000 customers make a purchase in the target window. A
model predicting no purchase for everyone achieves 95\% accuracy while
identifying no purchasers. A different model flags 100 customers,
including 30 purchasers and 70 nonpurchasers. It has 30 true positives,
20 false negatives, 70 false positives, and 880 true negatives.

The second model's precision is 30\%, recall is 60\%, and accuracy is
91\%. Its lower accuracy coexists with a potentially useful ranking for
a limited-capacity task. Whether it improves profit depends on the
action, costs, and whether the purchase probabilities under that action
are the relevant quantities. Accuracy alone cannot make that decision.

ROC curves compare true-positive and false-positive rates across
thresholds. Precision-recall curves focus on the tradeoff between
finding positives and the composition of the flagged group, which can be
informative when positives are rare. Both depend on the evaluation
setting; precision is especially sensitive to prevalence. Lift at a
specified contact fraction compares the selected group's outcome rate
with the population rate, but observed lift is not automatically an
incremental treatment effect.

Report metrics with their operational units. For a service team that can
review only 100 cases per week, precision among the top 100 and the
number of important cases missed may be more relevant than a global
ranking metric. Use a holdout that reflects the weekly population and
inspect uncertainty in the metric when counts are small.

\hypertarget{M07-ranking-and-calibration-answer-different-questions}{%
\section{Ranking and Calibration Answer Different
Questions}\label{M07-ranking-and-calibration-answer-different-questions}}

\noindent A ranking model orders cases from more likely to less likely. A
calibrated probability model assigns probabilities that correspond to
observed frequencies in relevant groups. If customers assigned
probability 0.20 purchase about 20\% of the time in a suitable
evaluation population, that group is calibrated. Calibration can vary
across segments and over time even when overall performance looks
satisfactory.

A strictly increasing transformation of scores preserves their ranking.
Squaring probabilities between zero and one therefore leaves a
ranking-based AUC unchanged, apart from practical handling of ties or
finite precision, while generally changing calibration. A predicted
probability of 0.50 becomes 0.25. If the original 0.50 group purchases
half the time, the transformed probabilities understate its frequency.

\begin{figure}[htbp]
\centering
% Original synthetic teaching figure; see corresponding module.
\begin{tikzpicture}
\begin{axis}[xmin=0,xmax=1,ymin=0,ymax=1,xlabel={Reported purchase probability},ylabel={Implied observed purchase frequency},xtick={0,.25,.5,.75,1},ytick={0,.25,.5,.75,1},legend style={at={(.5,1.08)},anchor=south,legend columns=2},ymajorgrids=true]
\addplot[IlliniBlue,dashed,domain=0:1] {x};
\addplot[IlliniOrange,very thick,domain=0:1,samples=151] {sqrt(x)};
\legend{Calibrated score,Squared score}
\addplot[IlliniBlue,only marks,mark=*] coordinates {(.25,.5)};
\node[align=left,anchor=north west,font=\scriptsize] at (axis cs:.48,.25) {The original 0.50 becomes 0.25.\\Ranking is unchanged.};
\end{axis}
\end{tikzpicture}

\caption{A monotonic score transformation preserves ranking while changing probability calibration.}\label{fig:M07_calibration}
\end{figure}

For calibrated forecasts, predicted probabilities equal observed
frequencies. After squaring an initially calibrated score \(p\), the
reported probability is \(q=p^2\), while the observed frequency remains
\(p=\sqrt{q}\). Expected counts and profit calculations depend on the
probability scale, so unchanged ranking performance does not guarantee
unchanged decision quality.

Inspect calibration plots and proper scoring rules, such as log loss or
the Brier score, on held-out data. Calibration methods themselves
require fitting and should use data separate from the final evaluation.
If the deployment population changes, recalibration may be necessary. A
single global correction may be insufficient when the relationship
changes differently across groups.

\hypertarget{M07-thresholds-should-reflect-the-decision}{%
\section{Thresholds Should Reflect the
Decision}\label{M07-thresholds-should-reflect-the-decision}}

\noindent Consider a classification action in which a false positive costs
\(C_{FP}\) and a false negative costs \(C_{FN}\), with zero cost for
correct classifications. For a calibrated probability \(p\), predicting
positive has expected cost \(C_{FP}(1-p)\) and predicting negative has
expected cost \(C_{FN}p\). The cost-minimizing rule predicts positive
when

\[
p>\frac{C_{FP}}{C_{FP}+C_{FN}}.
\]

If an unnecessary review costs one unit and missing a true issue costs
nine units, the threshold is 0.10 under this simplified loss. A
threshold of 0.50 would encode a different cost balance. Capacity
constraints, review effectiveness, and heterogeneous costs require a
more detailed rule. The formula is useful because it makes the decision
assumptions explicit.

A marketing intervention usually needs a causal comparison beyond this
classification loss. For a coupon, let \(p_i(1)\) and \(p_i(0)\) be
purchase probabilities with and without the coupon, \(m\) the
contribution before the discount, \(d\) the discount paid on a treated
purchase, and \(c\) the contact cost. A simple incremental expected
contribution is

\[
\Delta\pi_i=m\{p_i(1)-p_i(0)\}-d\,p_i(1)-c.
\]

The discount is paid on all treated purchases in this setup, including
purchases that would have happened anyway. The equation assumes one
relevant purchase opportunity and a fixed discount and contribution. It
can be extended to spending, multiple purchases, returns, and
heterogeneous margins.

\hypertarget{M07-case-study-targeting-the-most-responsive-customers-can-lose-money}{%
\section{Case Study: Targeting the Most Responsive Customers Can Lose
Money}\label{M07-case-study-targeting-the-most-responsive-customers-can-lose-money}}

\noindent Take \(m=40\), \(d=5\), and \(c=2\). Customer A has purchase
probabilities 0.80 with the coupon and 0.78 without it. Customer B has
probabilities 0.30 and 0.15. A has the higher treated response
probability, but incremental contribution is \$-5.20 for A and \$2.50
for B.

\begin{table}[H]
\centering\small
\begin{tabular}{@{}
  >{\raggedright\arraybackslash}p{(\columnwidth - 6\tabcolsep) * \real{0.20}}
  >{\raggedleft\arraybackslash}p{(\columnwidth - 6\tabcolsep) * \real{0.26}}
  >{\raggedleft\arraybackslash}p{(\columnwidth - 6\tabcolsep) * \real{0.26}}
  >{\raggedleft\arraybackslash}p{(\columnwidth - 6\tabcolsep) * \real{0.28}}@{}}
\toprule
Customer & Purchase With Coupon & Purchase Without Coupon & Incremental
Contribution \\
\midrule

A & 0.80 & 0.78 & -\$5.20 \\
B & 0.30 & 0.15 & \$2.50 \\
\bottomrule
\end{tabular}
\end{table}

\begin{figure}[htbp]
\centering
% Original synthetic teaching figure; see corresponding module.
\begin{tikzpicture}
\begin{axis}[ybar,bar width=15pt,xmin=.4,xmax=2.6,ymin=-7.5,ymax=8.5,xtick={1,2},xticklabels={Customer A,Customer B},ylabel={Contribution per customer (dollars)},ymajorgrids=true,legend style={at={(.5,1.13)},anchor=south,legend columns=3},nodes near coords,nodes near coords style={font=\scriptsize}]
\addplot[fill=IlliniBlue,draw=IlliniBlue] coordinates {(1,.8)(2,6)};
\addplot[fill=LightStorm,draw=Storm] coordinates {(1,-6)(2,-3.5)};
\addplot[fill=IlliniOrange,draw=IlliniOrange] coordinates {(1,-5.2)(2,2.5)};
\legend{Incremental purchases,Discount plus contact,Net contribution}
\end{axis}
\end{tikzpicture}

\caption{Baseline purchasing and discount expense can reverse a response-based targeting rule.}\label{fig:M07_targeting}
\end{figure}

Incremental purchasing contributes \$0.80 for A and \$6.00 for B, while
discount and contact costs are \$6.00 and \$3.50, respectively. These
values follow from the assumed purchase probabilities. Estimating their
difference from customer data requires treatment variation and an
identification strategy. Uplift or heterogeneous-treatment-effect models
address that causal target under their conditions; they are not ordinary
response models with a different name.

Evaluate a targeting policy on data that support its value comparison. A
randomized holdout or a carefully designed policy evaluation can
estimate incremental contribution for the selected group. Estimating a
model and reporting its value on the same observations used to choose
the policy can overstate performance. Budget constraints also matter: a
customer with positive expected value may still be excluded when another
customer offers greater value per unit of scarce capacity.

\hypertarget{M07-interpretation-fairness-and-deployment}{%
\section{Interpretation, Fairness, and
Deployment}\label{M07-interpretation-fairness-and-deployment}}

\noindent Feature importance describes a relationship between a fitted model and
its inputs under a chosen method. Permutation importance measures the
change in predictive performance after disrupting a feature, with
complications when predictors are correlated. Partial-dependence plots
average predictions over a reference distribution and can involve
unrealistic combinations of correlated features. Local explanation
methods describe the fitted prediction under their conventions. These
tools do not automatically identify causal effects of changing a
feature.

Evaluate performance across groups relevant to the application and
examine how errors translate into service or opportunity. A model can be
accurate on average while performing poorly for new customers, small
stores, or a language group with sparse training data. Protected or
sensitive attributes and proxies require appropriate institutional and
legal review for the intended use. The consequences of errors depend on
the decision and the affected customers.

Deployment requires monitoring the input distribution, missingness,
outcome frequency, calibration, and decision value. A change in model
accuracy can reflect changing customer behavior, instrumentation, or
policy. Record versions and retain a reproducible baseline. If an
intervention changes the behavior that generates future training data,
the model's environment becomes endogenous to its use, because its
recommendations change subsequent exposure and behavior.

\hypertarget{M07-ai-assisted-workflow-practice-and-readings}{%
\section{AI-Assisted Workflow, Practice, and
Readings}\label{M07-ai-assisted-workflow-practice-and-readings}}

\noindent Ask an assistant to create an evaluation plan before choosing a model.
Specify forecast origin, horizon, unit, grouping, baseline, loss, and
final holdout. Require a feature-availability table and a pipeline that
fits transformations only on training data. Have it produce both a
metric table and an explanation of the operational decision each metric
supports.

For practice, verify the 91\% accuracy and 60\% recall in the
rare-purchase example. Square a small set of predicted probabilities and
confirm that their order is unchanged while expected purchase counts
change. Recalculate the coupon case with a \$1 discount and explain
which customer becomes more attractive. Finally, identify why a model
trained on all customers' latest outcomes cannot honestly claim to
predict those same outcomes out of sample.

\href{https://www.statlearning.com/}{An Introduction to Statistical
Learning} develops the central learning methods. The
\href{https://scikit-learn.org/stable/modules/model_evaluation.html}{scikit-learn
evaluation guide} provides metric definitions and implementation
references. \href{https://doi.org/10.1007/s11002-020-09535-7}{Hagen and
colleagues (2020)} connect machine-learning tools to behavioral
marketing questions.

\hypertarget{M08-m8.-current-topics-and-integrated-marketing-decisions}{%
\chapter{Current Topics and Integrated Marketing
Decisions}\label{M08-m8.-current-topics-and-integrated-marketing-decisions}}

\noindent A marketing team can now generate campaign copy, build a prediction
model, and summarize customer feedback within a short working session.
The remaining difficulty is deciding whether the resulting system
improves a business decision. A recommendation engine can reproduce past
exposure, a media model can attribute seasonal demand to advertising,
and an AI-generated analysis can calculate an invalid comparison
precisely. Evaluating these systems requires a clear account of their
information, incentives, assumptions, and consequences.

\hypertarget{M08-recommendation-systems-learn-from-selected-exposure}{%
\section{Recommendation Systems Learn From Selected
Exposure}\label{M08-recommendation-systems-learn-from-selected-exposure}}

\noindent A retailer wants to recommend products that customers will find useful.
A common starting point is past interaction: purchases, clicks, ratings,
or time spent. These observations combine preferences with the
opportunities customers received. A product displayed prominently has
more chances to be clicked than an otherwise similar product that was
rarely shown. An absent interaction can therefore mean dislike, lack of
exposure, lack of need, or an unavailable product.

Content-based methods use product attributes and customer histories to
identify similar items. Collaborative methods use patterns across users
and items, including matrix factorization, which approximates an
interaction matrix with lower-dimensional user and item representations.
Hybrid systems combine these sources. Each method needs a strategy for
new customers and new products that have little interaction history.
Popularity baselines and simple category rules provide useful
comparisons.

Recommendation systems also change the data used to train the next
model. A product that receives more exposure can accumulate more
interactions, become easier to estimate, and receive still more
exposure. This process can amplify initial differences without
establishing that the repeatedly displayed product would be preferred
under equal exposure.

\begin{figure}[htbp]
\centering
% Original synthetic teaching figure; see corresponding module.
\begin{tikzpicture}
\node[box,text width=3cm] (m) at (0,0) {Model and ranking};
\node[box,text width=3cm] (e) at (5,0) {Product exposure};
\node[box,text width=3cm] (y) at (10,0) {Observed behavior};
\node[orangebox,text width=3cm] (p) at (10,2.4) {Customer preferences};
\node[box,text width=3cm] (a) at (5,2.4) {Availability and position};
\draw[edge] (m)--(e);\draw[edge] (e)--(y);\draw[edge] (p)--(y);\draw[edge] (a)--(e);
\draw[accent] (y.south)--(10,-1.5)--node[below] {New training data} (0,-1.5)--(m.south);
\node[align=center,text width=11cm] at (5,-2.7) {Interaction data reflect both preference and opportunity.\\A feedback loop can reinforce initially favored items.};
\end{tikzpicture}

\caption{A recommender changes the exposure that generates its next training sample.}\label{fig:M08_feedback}
\end{figure}

Rankings determine exposure, exposure affects observed behavior, and
that behavior enters subsequent training data. Customer preferences also
affect behavior, while availability and position affect the opportunity
to respond. Logging impressions, positions, eligible alternatives, and
the policy that selected them helps distinguish these components.
Purchase logs alone often omit the alternatives a customer could
actually have encountered.

Evaluate recommendation quality at the unit relevant to the decision.
Precision among the top recommended items and ranking metrics assess
aspects of relevance. Coverage, novelty, diversity, repeat engagement,
returns, and contribution describe other consequences. A system can
improve one metric while reducing another. Offline evaluation on
historical logs is particularly sensitive to the policy that generated
those logs. An experiment that changes exposure can supply information
that a more flexible model of the same selected data cannot supply.

\hypertarget{M08-case-study-relevance-variety-and-customer-goals}{%
\section{Case Study: Relevance, Variety, and Customer
Goals}\label{M08-case-study-relevance-variety-and-customer-goals}}

\noindent Nielsen, Yang, Dubé, Knäuper, Ling, and Nie study consumption variety in
food recommendation. Their research connects observed patterns in food
logs with a recommendation method that incorporates variety. Ling and
coauthors develop a goal-oriented recommendation approach using behavior
from comparable users. The choice of objective matters: reproducing
observed choices can conflict with a customer's stated goal.
Observational relationships and offline ranking evaluations measure
different objects from the causal effect of deploying a recommendation
policy.

Suppose a retail customer requests three products and values variety
across categories. Assign four candidate items the relevance scores
below. The scores are comparable across these items and measure
relevance rather than purchase probability.

\begin{table}[H]
\centering\small
\begin{tabular}{@{}llr@{}}
\toprule
Item & Category & Relevance Score \\
\midrule

A & Coffee & 0.90 \\
B & Coffee & 0.85 \\
C & Tea & 0.75 \\
D & Fruit & 0.65 \\
\bottomrule
\end{tabular}
\end{table}

Choosing the three highest scores gives A, B, and C, with total
relevance 2.50 and two categories. Choosing A, C, and D gives relevance
2.30 and three categories. To make the tradeoff explicit, define a list
objective

\[
J(L)=\sum_{j\in L}s_j+\lambda K(L),
\]

where \(L\) is a three-item list, \(s_j\) is relevance, \(K(L)\) counts
distinct categories, and \(\lambda\) is the value placed on an
additional category in the same units as the relevance score. With
\(\lambda=0.30\), A, B, C has objective 3.10, while A, C, D has
objective 3.20. The latter list is preferred under this stated
objective. The threshold at which these two lists tie is
\(\lambda=0.20\).

The variety weight must come from an explicit objective, user
preference, or evaluation process. Category counts are also a coarse
measure: they treat all category differences equally and ignore which
products the customer already owns. A deployed system would need better
measurements and an evaluation of whether the resulting lists improve
the intended customer experience.

For an experiment, randomize eligible customers to the existing and
revised recommendation policies. Choose a primary outcome such as
contribution per assigned customer over a fixed horizon, alongside
measures of customer satisfaction and unwanted repetition. Analyze
assignment rather than only customers who clicked a recommendation.
Clicking occurs after the policy is assigned and can differ between
policies. Preserve the list shown and the alternatives eligible at each
impression so the mechanism can be examined.

\hypertarget{M08-marketing-mix-models-combine-carryover-and-saturation}{%
\section{Marketing-Mix Models Combine Carryover and
Saturation}\label{M08-marketing-mix-models-combine-carryover-and-saturation}}

\noindent A marketing-mix model relates aggregate outcomes to media activity and
other determinants across time or markets. Its appeal is practical: many
firms observe weekly sales and channel spending even when individual
customer paths are unavailable. The analytical challenge is that media
decisions often respond to expected demand. Advertising, promotions,
holidays, distribution, and competitor actions can move together.

A simple carryover variable is

\[
a_{ct}=s_{ct}+\rho_c a_{c,t-1},\qquad 0\leq\rho_c<1,
\]

where \(s_{ct}\) is spending in channel \(c\) during period \(t\),
\(a_{ct}\) is its accumulated exposure proxy, and \(\rho_c\) controls
geometric decay. This unnormalized definition gives the same units as
spending, but its long-run level under constant spending is
\(s_c/(1-\rho_c)\). Some software normalizes the weights. Coefficients
and interpretations must follow the selected convention.

If one spending pulse of 100 occurs at period zero and \(\rho=0.5\), the
subsequent accumulated values are 100, 50, 25, and 12.5. A pulse can
therefore coincide with later sales even when current spending is zero.
A seasonal demand pattern can also create persistence. The model needs
enough information and a credible design to distinguish these
explanations.

Saturation describes declining additional response at higher exposure. A
saturating response curve is \(R(a)=A[1-\exp(-a/b)]\), where \(A\) is
the limiting response and \(b>0\) controls how quickly the curve
approaches that limit. A basic outcome model may combine channel
response curves, controls, trend, seasonality, and an error term. The
fitted components are model-dependent; their causal interpretation
requires assumptions about how spending was chosen and which demand
determinants were measured.

The official
\href{https://developers.google.com/meridian/docs/basics/meridian-introduction}{Meridian
documentation} describes a modern framework for marketing-mix modeling.
The
\href{https://www.pymc-marketing.io/en/latest/notebooks/mmm/mmm_example.html}{PyMC
Marketing examples} provide another implementation route. Use those
materials to inspect definitions, priors, and diagnostics. A software
package cannot supply a missing source of identifying variation.
Experiments and external evidence can help calibrate a model, subject to
whether their populations, interventions, and time horizons match the
intended use.

\hypertarget{M08-case-study-allocate-the-next-dollar-using-marginal-returns}{%
\section{Case Study: Allocate the Next Dollar Using Marginal
Returns}\label{M08-case-study-allocate-the-next-dollar-using-marginal-returns}}

\noindent Consider two channels with imposed contribution response curves

\[
R_A(s_A)=100[1-\exp(-s_A/20)],\qquad
R_B(s_B)=80[1-\exp(-s_B/40)].
\]

Both spending and contribution are measured in thousands of dollars over
the same horizon. These response functions omit cross-channel
interactions and are assumed known. At a total budget of 60, suppose the
current allocation is 40 to A and 20 to B. The question is whether
moving a small amount between channels increases contribution while
holding the total budget fixed.

Average contribution per unit spent is \(R_c(s_c)/s_c\). At the current
allocation, it is approximately 2.16 for A and 1.57 for B. Marginal
contribution is the derivative of the response curve. It is
approximately 0.68 for A and 1.21 for B. A has the higher average
return, while B offers the greater return on the next small unit of
spending.

\begin{figure}[htbp]
\centering
% Original synthetic teaching figure; see corresponding module.
\begin{tikzpicture}
\begin{axis}[xmin=0,xmax=65,ymin=0,ymax=110,xlabel={Channel spending (thousands of dollars)},ylabel={Contribution (thousands of dollars)},legend style={at={(.5,1.08)},anchor=south,legend columns=2},ymajorgrids=true]
\addplot[IlliniBlue,very thick,domain=0:60,samples=120] {100*(1-exp(-x/20))};
\addplot[IlliniOrange,very thick,domain=0:60,samples=120] {80*(1-exp(-x/40))};
\legend{Channel A,Channel B}
\addplot[IlliniBlue,dashed,domain=30:50] {86.4665+.676676*(x-40)};
\addplot[IlliniOrange,dashed,domain=10:30] {31.4775+1.21306*(x-20)};
\addplot[only marks,mark=*,IlliniBlue,mark size=3pt] coordinates {(40,86.4665)};
\addplot[only marks,mark=square*,IlliniOrange,mark size=3pt] coordinates {(20,31.4775)};
\node[anchor=south east,font=\scriptsize] at (axis cs:40,93) {At 40: slope 0.68};
\node[anchor=north west,font=\scriptsize] at (axis cs:22,28) {At 20: slope 1.21};
\end{axis}
\end{tikzpicture}

\caption{Average returns can favor one channel while marginal returns favor the other.}\label{fig:M08_marginal_returns}
\end{figure}

At the current allocations, B has the steeper response curve. Moving a
small amount from A to B therefore raises total contribution. The
comparison depends on the slopes at the current allocation, whereas
average returns summarize all spending up to that point. The budget is
held fixed, so the transfer does not change total spending expense.

For a differentiable, concave response with an interior optimum and no
other constraints, marginal returns are equal across channels. Here the
condition is

\[
5\exp(-s_A/20)=2\exp(-s_B/40),\qquad s_A+s_B=60.
\]

Solving gives \(s_A\approx32.22\) and \(s_B\approx27.78\). Total
contribution rises from approximately 117.94 to 120.09, or about 2.14
thousand dollars. If total budget were also a choice, marginal
contribution would be compared with the marginal cost of spending;
exhausting a fixed budget and choosing its size are different
optimization problems.

A real decision should also examine parameter uncertainty, operational
minimums, saturation outside observed spending ranges, delayed effects,
and interactions across channels. Report how the recommended allocation
changes under plausible response curves. A sharply stated optimum based
on weakly identified slopes can conceal substantial decision
uncertainty. Prefer an allocation that is supported by observed
variation and specify an experiment that would resolve the most
consequential uncertainty.

\hypertarget{M08-generative-ai-changes-the-cost-of-producing-and-checking-work}{%
\section{Generative AI Changes the Cost of Producing and Checking
Work}\label{M08-generative-ai-changes-the-cost-of-producing-and-checking-work}}

\noindent Generative AI can assist with code, text classification, documentation,
and communication. The business value depends on the complete workflow.
Faster generation can be offset by additional review, integration
failures, or errors that reach downstream decisions. Measure the output
that matters: a correct, reproducible analysis completed with an
appropriate understanding of its limitations.

Separate several uses. A coding assistant produces or edits a program
that can be inspected and executed. An annotation model produces labels
that require a measurement evaluation. A retrieval system selects source
passages and may generate a summary that needs factual support. A
content system produces customer-facing material whose effect should be
evaluated in its deployment context. Each use has a different unit of
quality and a different failure mechanism.

For coding, begin with an analytical specification and a small example
with a known result. Ask the assistant to implement a limited step,
explain the data transformations, and identify assumptions. Inspect the
code before running it. Then compare intermediate outputs with
independent checks. A program that executes successfully can still use
an incorrect denominator, join on the wrong key, or include future
information. Those errors require substantive checks rather than syntax
repair.

For retrieval and summarization, maintain the connection between a claim
and its source. A citation should support the particular statement next
to it. Ask the model to distinguish an observed result from a proposed
explanation, and check that distinction against the original text. If a
passage or file is unavailable, state that limitation. Fluent
reconstruction is not a substitute for access to the evidence.

\hypertarget{M08-case-study-evaluate-an-ai-assisted-analytics-workflow}{%
\section{Case Study: Evaluate an AI-Assisted Analytics
Workflow}\label{M08-case-study-evaluate-an-ai-assisted-analytics-workflow}}

\noindent Suppose a firm wants to know whether access to a coding assistant
improves analysts' work on customer-retention tasks. Define treatment as
access under a specified protocol and version, and define an outcome
that incorporates both completion and validity. A useful primary outcome
could be the fraction of assigned tasks meeting a prespecified
analytical rubric within a fixed time budget. Secondary outcomes could
include completion time, independent error severity, and a later
explanation task completed without assistance.

Randomize analysts or teams according to the likely spillovers. Repeated
tasks by the same analyst require attention to learning, carryover, and
dependence. If a crossover design is used, randomize order and consider
whether experience in the first condition changes performance in the
second. A task-level comparison with no accounting for analyst ability
or task difficulty can confound the treatment effect with who received
which work.

Quality review should use a common rubric and, where feasible, reviewers
unaware of treatment assignment. The rubric can check observation units,
joins, leakage, estimation, uncertainty, and interpretation. Record the
assigned task, the eligible inputs, the tool configuration, and the
evaluation procedure. Avoid interpreting a favorable result for one task
type as a universal effect across all analytical work.

Suppose a manual workflow uses 90 minutes to produce an analysis and 20
minutes to review it, totaling 110 minutes. An assisted workflow uses 45
minutes to produce a draft and 40 minutes to review it, totaling 85
minutes. The reduction in total time is 25 minutes, or about 22.7\%,
while the drafting-time reduction is 50\%. The total-time comparison is
informative only if the completed outputs meet comparable quality
requirements.

\begin{figure}[htbp]
\centering
% Original synthetic teaching figure; see corresponding module.
\begin{tikzpicture}[x=.095cm,y=1cm]
\node[anchor=east] at (-2,1.5) {Manual};\node[anchor=east] at (-2,0) {AI-assisted};
\fill[IlliniBlue] (0,1.1) rectangle (90,1.9);\fill[IlliniOrange] (90,1.1) rectangle (110,1.9);
\fill[IlliniBlue] (0,-.4) rectangle (45,.4);\fill[IlliniOrange] (45,-.4) rectangle (85,.4);
\node[text=white] at (45,1.5) {Produce: 90};\node[text=IlliniBlue,font=\scriptsize,align=center] at (100,1.5) {Review:\\20};
\node[text=white] at (22.5,0) {Produce: 45};\node[text=IlliniBlue] at (65,0) {Review: 40};
\node[anchor=west] at (112,1.5) {110};\node[anchor=west] at (87,0) {85};
\draw[edge] (0,-1)--(120,-1) node[right] {Minutes};
\foreach \x in {0,30,60,90} {\draw[IlliniBlue] (\x,-1.05)--(\x,-.95);\node[anchor=north,font=\scriptsize] at (\x,-1.1) {\x};}
\node[align=center,text width=12cm] at (58,-2.2) {Drafting time falls 50\%; total time falls 22.7\%.\\Task and completed-work quality held fixed.};
\end{tikzpicture}

\caption{The relevant time saving includes generation, review, and correction.}\label{fig:M08_workflow_time}
\end{figure}

This comparison holds the task and required quality fixed. Review is
part of production because an undetected error can alter a marketing
decision. In empirical work, report the completion and quality outcomes
alongside time so that faster abandonment or lower-quality work is not
counted as an equivalent improvement.

\hypertarget{M08-integrated-case-a-retail-subscription-launch}{%
\section{Integrated Case: A Retail Subscription
Launch}\label{M08-integrated-case-a-retail-subscription-launch}}

\noindent A retailer is considering a subscription service with recommendations
and delivery benefits. Begin with the decision: which offer to launch,
to which eligible customers, at what price, with what evidence of
incremental contribution? Define the evaluation horizon, capacity
constraints, and the alternatives to launch.

Construct the customer and order tables with explicit keys and dates.
Describe customer activity and purchasing cohorts, keeping unavailable
future retention separate from observed inactivity. Use regression to
summarize conditional patterns while examining confounding. Design a
conjoint exercise to learn stated tradeoffs between delivery,
assortment, and price, then assess whether its predictions generalize to
an actual offer.

Use segmentation to organize meaningful differences in needs and
behavior, and calculate customer value from contribution rather than
revenue. Analyze reviews to identify the service attributes customers
mention, checking annotation quality against human labels. Develop a
prediction model using only information available at the targeting date.
Its forecasts can help plan capacity, while a randomized offer evaluates
incremental take-up, contribution, and retention.

Finally, evaluate the recommendation component and the acquisition
budget. Log exposure to recommended items, distinguish average from
marginal media returns, and use AI assistance within a documented
verification process. The management recommendation should identify the
supported action, the uncertainty that could change it, and the next
observation or experiment with the greatest decision value. A
complicated collection of models is useful only if these pieces support
a coherent choice.

\hypertarget{M08-practice-and-further-reading}{%
\section{Practice and Further
Reading}\label{M08-practice-and-further-reading}}

\noindent \textbf{Recommendation exercise.} Compare all four three-item lists in
the variety example at \(\lambda=0\), 0.20, and 0.30. At zero, A, B, C
is optimal. At 0.20, it ties A, C, D. At 0.30, A, C, D is optimal.
Explain how a different category definition could change that conclusion
without changing any relevance score.

\textbf{Budget exercise.} Calculate the gain from moving one thousand
dollars from channel A to channel B at the original allocation. Compare
the exact change with the difference between the two marginal returns.
The derivative is a local approximation; the exact finite transfer also
reflects curvature. Then explain why an unobserved seasonal demand shock
could invalidate a causal interpretation of an estimated response curve.

\textbf{AI evaluation exercise.} Draft a rubric for an
assistant-generated retention analysis. Include at least one check that
detects a substantive error in code that executes successfully, such as
an order-line denominator used for average order value. Define a
hand-checkable example that exposes the error and explain how a reviewer
would determine whether it occurred.

Read Nielsen and coauthors,
\href{https://doi.org/10.1086/720453}{``Consumption Variety in Food
Recommendation''}, and Ling and coauthors,
\href{https://doi.org/10.1145/3485447.3514193}{``Following Good
Examples''}, for applications connecting recommendation objectives with
customer behavior. The official marketing-mix documentation linked above
develops implementation details. Consult the
\href{../01_Resources/10_AI_Coding_Concepts_and_Evidence/README.md}{AI-coding
evidence} and
\href{../01_Resources/11_AI_Coding_Tools_and_Workflows/README.md}{workflow
resources} for further study.

\end{document}
